% Options for packages loaded elsewhere
% Options for packages loaded elsewhere
\PassOptionsToPackage{unicode}{hyperref}
\PassOptionsToPackage{hyphens}{url}
\PassOptionsToPackage{dvipsnames,svgnames,x11names}{xcolor}
%
\documentclass[
  letterpaper,
  DIV=11,
  numbers=noendperiod]{scrreprt}
\usepackage{xcolor}
\usepackage{amsmath,amssymb}
\setcounter{secnumdepth}{5}
\usepackage{iftex}
\ifPDFTeX
  \usepackage[T1]{fontenc}
  \usepackage[utf8]{inputenc}
  \usepackage{textcomp} % provide euro and other symbols
\else % if luatex or xetex
  \usepackage{unicode-math} % this also loads fontspec
  \defaultfontfeatures{Scale=MatchLowercase}
  \defaultfontfeatures[\rmfamily]{Ligatures=TeX,Scale=1}
\fi
\usepackage{lmodern}
\ifPDFTeX\else
  % xetex/luatex font selection
\fi
% Use upquote if available, for straight quotes in verbatim environments
\IfFileExists{upquote.sty}{\usepackage{upquote}}{}
\IfFileExists{microtype.sty}{% use microtype if available
  \usepackage[]{microtype}
  \UseMicrotypeSet[protrusion]{basicmath} % disable protrusion for tt fonts
}{}
\makeatletter
\@ifundefined{KOMAClassName}{% if non-KOMA class
  \IfFileExists{parskip.sty}{%
    \usepackage{parskip}
  }{% else
    \setlength{\parindent}{0pt}
    \setlength{\parskip}{6pt plus 2pt minus 1pt}}
}{% if KOMA class
  \KOMAoptions{parskip=half}}
\makeatother
% Make \paragraph and \subparagraph free-standing
\makeatletter
\ifx\paragraph\undefined\else
  \let\oldparagraph\paragraph
  \renewcommand{\paragraph}{
    \@ifstar
      \xxxParagraphStar
      \xxxParagraphNoStar
  }
  \newcommand{\xxxParagraphStar}[1]{\oldparagraph*{#1}\mbox{}}
  \newcommand{\xxxParagraphNoStar}[1]{\oldparagraph{#1}\mbox{}}
\fi
\ifx\subparagraph\undefined\else
  \let\oldsubparagraph\subparagraph
  \renewcommand{\subparagraph}{
    \@ifstar
      \xxxSubParagraphStar
      \xxxSubParagraphNoStar
  }
  \newcommand{\xxxSubParagraphStar}[1]{\oldsubparagraph*{#1}\mbox{}}
  \newcommand{\xxxSubParagraphNoStar}[1]{\oldsubparagraph{#1}\mbox{}}
\fi
\makeatother


\usepackage{longtable,booktabs,array}
\usepackage{calc} % for calculating minipage widths
% Correct order of tables after \paragraph or \subparagraph
\usepackage{etoolbox}
\makeatletter
\patchcmd\longtable{\par}{\if@noskipsec\mbox{}\fi\par}{}{}
\makeatother
% Allow footnotes in longtable head/foot
\IfFileExists{footnotehyper.sty}{\usepackage{footnotehyper}}{\usepackage{footnote}}
\makesavenoteenv{longtable}
\usepackage{graphicx}
\makeatletter
\newsavebox\pandoc@box
\newcommand*\pandocbounded[1]{% scales image to fit in text height/width
  \sbox\pandoc@box{#1}%
  \Gscale@div\@tempa{\textheight}{\dimexpr\ht\pandoc@box+\dp\pandoc@box\relax}%
  \Gscale@div\@tempb{\linewidth}{\wd\pandoc@box}%
  \ifdim\@tempb\p@<\@tempa\p@\let\@tempa\@tempb\fi% select the smaller of both
  \ifdim\@tempa\p@<\p@\scalebox{\@tempa}{\usebox\pandoc@box}%
  \else\usebox{\pandoc@box}%
  \fi%
}
% Set default figure placement to htbp
\def\fps@figure{htbp}
\makeatother


% definitions for citeproc citations
\NewDocumentCommand\citeproctext{}{}
\NewDocumentCommand\citeproc{mm}{%
  \begingroup\def\citeproctext{#2}\cite{#1}\endgroup}
\makeatletter
 % allow citations to break across lines
 \let\@cite@ofmt\@firstofone
 % avoid brackets around text for \cite:
 \def\@biblabel#1{}
 \def\@cite#1#2{{#1\if@tempswa , #2\fi}}
\makeatother
\newlength{\cslhangindent}
\setlength{\cslhangindent}{1.5em}
\newlength{\csllabelwidth}
\setlength{\csllabelwidth}{3em}
\newenvironment{CSLReferences}[2] % #1 hanging-indent, #2 entry-spacing
 {\begin{list}{}{%
  \setlength{\itemindent}{0pt}
  \setlength{\leftmargin}{0pt}
  \setlength{\parsep}{0pt}
  % turn on hanging indent if param 1 is 1
  \ifodd #1
   \setlength{\leftmargin}{\cslhangindent}
   \setlength{\itemindent}{-1\cslhangindent}
  \fi
  % set entry spacing
  \setlength{\itemsep}{#2\baselineskip}}}
 {\end{list}}
\usepackage{calc}
\newcommand{\CSLBlock}[1]{\hfill\break\parbox[t]{\linewidth}{\strut\ignorespaces#1\strut}}
\newcommand{\CSLLeftMargin}[1]{\parbox[t]{\csllabelwidth}{\strut#1\strut}}
\newcommand{\CSLRightInline}[1]{\parbox[t]{\linewidth - \csllabelwidth}{\strut#1\strut}}
\newcommand{\CSLIndent}[1]{\hspace{\cslhangindent}#1}



\setlength{\emergencystretch}{3em} % prevent overfull lines

\providecommand{\tightlist}{%
  \setlength{\itemsep}{0pt}\setlength{\parskip}{0pt}}



 


\KOMAoption{captions}{tableheading}
\makeatletter
\@ifpackageloaded{tcolorbox}{}{\usepackage[skins,breakable]{tcolorbox}}
\@ifpackageloaded{fontawesome5}{}{\usepackage{fontawesome5}}
\definecolor{quarto-callout-color}{HTML}{909090}
\definecolor{quarto-callout-note-color}{HTML}{0758E5}
\definecolor{quarto-callout-important-color}{HTML}{CC1914}
\definecolor{quarto-callout-warning-color}{HTML}{EB9113}
\definecolor{quarto-callout-tip-color}{HTML}{00A047}
\definecolor{quarto-callout-caution-color}{HTML}{FC5300}
\definecolor{quarto-callout-color-frame}{HTML}{acacac}
\definecolor{quarto-callout-note-color-frame}{HTML}{4582ec}
\definecolor{quarto-callout-important-color-frame}{HTML}{d9534f}
\definecolor{quarto-callout-warning-color-frame}{HTML}{f0ad4e}
\definecolor{quarto-callout-tip-color-frame}{HTML}{02b875}
\definecolor{quarto-callout-caution-color-frame}{HTML}{fd7e14}
\makeatother
\makeatletter
\@ifpackageloaded{bookmark}{}{\usepackage{bookmark}}
\makeatother
\makeatletter
\@ifpackageloaded{caption}{}{\usepackage{caption}}
\AtBeginDocument{%
\ifdefined\contentsname
  \renewcommand*\contentsname{Table of contents}
\else
  \newcommand\contentsname{Table of contents}
\fi
\ifdefined\listfigurename
  \renewcommand*\listfigurename{List of Figures}
\else
  \newcommand\listfigurename{List of Figures}
\fi
\ifdefined\listtablename
  \renewcommand*\listtablename{List of Tables}
\else
  \newcommand\listtablename{List of Tables}
\fi
\ifdefined\figurename
  \renewcommand*\figurename{Figure}
\else
  \newcommand\figurename{Figure}
\fi
\ifdefined\tablename
  \renewcommand*\tablename{Table}
\else
  \newcommand\tablename{Table}
\fi
}
\@ifpackageloaded{float}{}{\usepackage{float}}
\floatstyle{ruled}
\@ifundefined{c@chapter}{\newfloat{codelisting}{h}{lop}}{\newfloat{codelisting}{h}{lop}[chapter]}
\floatname{codelisting}{Listing}
\newcommand*\listoflistings{\listof{codelisting}{List of Listings}}
\makeatother
\makeatletter
\makeatother
\makeatletter
\@ifpackageloaded{caption}{}{\usepackage{caption}}
\@ifpackageloaded{subcaption}{}{\usepackage{subcaption}}
\makeatother
\usepackage{bookmark}
\IfFileExists{xurl.sty}{\usepackage{xurl}}{} % add URL line breaks if available
\urlstyle{same}
\hypersetup{
  pdftitle={Unjust Religion as Fiction},
  pdfauthor={Montaque Reynolds},
  colorlinks=true,
  linkcolor={blue},
  filecolor={Maroon},
  citecolor={Blue},
  urlcolor={Blue},
  pdfcreator={LaTeX via pandoc}}


\title{Unjust Religion as Fiction}
\author{Montaque Reynolds}
\date{2026-03-09}
\begin{document}
\maketitle

\renewcommand*\contentsname{Table of contents}
{
\hypersetup{linkcolor=}
\setcounter{tocdepth}{2}
\tableofcontents
}

\bookmarksetup{startatroot}

\chapter{Working Papers 2026 -- Status
Dashboard}\label{working-papers-2026-status-dashboard}

Current status as of
\texttt{r\ format(Sys.Date(),\ \textquotesingle{}\%B\ \%Y\textquotesingle{})}

\section{Active Papers}\label{active-papers}

\begin{tcolorbox}[enhanced jigsaw, bottomrule=.15mm, colback=white, breakable, arc=.35mm, rightrule=.15mm, toprule=.15mm, left=2mm, leftrule=.75mm, opacityback=0, colframe=quarto-callout-important-color-frame]

\vspace{-3mm}\textbf{Manticore}\vspace{3mm}

\begin{itemize}
\tightlist
\item[$\square$]
  Literature Review
\item[$\square$]
  Primary texts read \& noted
\item[$\square$]
  Brainstorm Core Argument
\item[$\square$]
  Draft paper outline with main sections
\item[$\square$]
  Write first draft
\item[$\square$]
  Share draft with trusted colleagues
\item[$\square$]
  Identify possible objections
\item[$\square$]
  Research Relevant Conferences
\item[$\square$]
  Research target journals for submission
\item[$\square$]
  Objections section started
\item[$\square$]
  Responses to objections
\item[$\square$]
  Submit to journal X by end of Q2
\end{itemize}

\end{tcolorbox}

\begin{tcolorbox}[enhanced jigsaw, bottomrule=.15mm, colback=white, breakable, arc=.35mm, rightrule=.15mm, toprule=.15mm, left=2mm, leftrule=.75mm, opacityback=0, colframe=quarto-callout-tip-color-frame]

\vspace{-3mm}\textbf{The Aesthetic Problem of Evil}\vspace{3mm}

\begin{itemize}
\tightlist
\item[$\square$]
  Literature Review
\item[$\square$]
  Primary texts read \& noted
\item[$\square$]
  Brainstorm Core Argument
\item[$\square$]
  Draft paper outline with main sections
\item[$\square$]
  Write first draft
\item[$\square$]
  Share draft with trusted colleagues
\item[$\square$]
  Identify possible objections
\item[$\square$]
  Research Relevant Conferences
\item[$\square$]
  Research target journals for submission
\item[$\square$]
  Objections section started
\item[$\square$]
  Responses to objections
\item[$\square$]
  Submit to journal X by end of Q2
\end{itemize}

\subsection{Morality as Privilege}\label{morality-as-privilege}

\begin{itemize}
\tightlist
\item[$\square$]
  Literature Review
\item[$\square$]
  Primary texts read \& noted
\item[$\square$]
  Brainstorm Core Argument
\item[$\square$]
  Draft paper outline with main sections
\item[$\square$]
  Write first draft
\item[$\square$]
  Share draft with trusted colleagues
\item[$\square$]
  Identify possible objections
\item[$\square$]
  Research Relevant Conferences
\item[$\square$]
  Research target journals for submission
\item[$\square$]
  Objections section started
\item[$\square$]
  Responses to objections
\item[$\square$]
  Submit to journal X by end of Q2
\end{itemize}

\subsection{God as the Best Possible Epistemic Ground of Human
Dignity}\label{god-as-the-best-possible-epistemic-ground-of-human-dignity}

\begin{itemize}
\tightlist
\item[$\square$]
  Literature Review

  \begin{itemize}
  \tightlist
  \item[$\square$]
    Imago Dei as Human Identity by Ryan Peterson
  \end{itemize}
\item[$\square$]
  Primary texts read \& noted
\item[$\square$]
  Brainstorm Core Argument
\item[$\square$]
  Draft paper outline with main sections
\item[$\square$]
  Write first draft
\item[$\square$]
  Share draft with trusted colleagues
\item[$\square$]
  Identify possible objections
\item[$\square$]
  Research Relevant Conferences
\item[$\square$]
  Research target journals for submission
\item[$\square$]
  Objections section started
\item[$\square$]
  Responses to objections
\item[$\square$]
  Submit to journal X by end of Q2
\end{itemize}

\end{tcolorbox}

\section{Next Up}\label{next-up}

\begin{itemize}
\tightlist
\item
  \href{./papers/paper-1-aes-evi/index.qmd}{Aesthetics of Evil}
\end{itemize}

\part{Active Papers}

\chapter{Aesthetics of Evil}\label{aesthetics-of-evil}

Montaque Reynolds (Stetson University)

\hfill\break

\section{Abstract}\label{abstract}

In medieval bestiaries, there are two classes of evil represented by
beasts; those which deceive and those which generate fear and death.
Among the first class of beasts, those which are notable for their
abilities to deceive, there are three kinds of deception: these are
noted by the camouflages of the chameleon, accurate imitations of the
manticore, and the ``seductive skills of sirens''.

\begin{itemize}
\tightlist
\item
  \textbf{Deceivers}: Camouflage, imitation, seduction

  \begin{itemize}
  \tightlist
  \item
    Chameleon: camouflage
  \item
    Manticore: accurate imitations
  \item
    Sirens: seductive skills
  \end{itemize}
\item
  \textbf{Fear and Death Generators}
\end{itemize}

While we might think this to be a medieval taxonomy of epistemic harms
only, the kinds of epistemic harms the medievalists worried about are
not far off from contemporary epistemic harms and risks; misinformation,
false and inauthentic narratives, and bullsh*t. This article looks at
those beasts which represent epistemic harms and attempts to develop an
account of inauthentic narrative as a moral harm.

On some accounts, the manticore is a union between a hyena and Ethiopian
lioness. It imitates the voices of men and possesses a fixed and
immoveable gaze. In addition, it does not possess gums in either of its
jaws, and rather than having multiple teeth, it has one large one that
curves around its jaw bone.

Pliny the elder is one who describes the manticore as a union between an
Ethiopian lioness and hyaena. Two notoriously deadly enemies who fight
for the same for supplies across the African plain. The fact that he
draws emphasis on the fact that the hyaena is male while the lioness is
female suggests that the manticore is born from a procreative union
between the two, and not merely a creative one. Perhaps the manticore
has obtained its ability to speak from the male hyaena's ability to
imitate the voice of a man. Pliny tells us that the hyaena is able to
learn a man's name, so that in calling him away from the protection of
his home, can then devour him. It appears that the manticore too
possesses this ability, and not only of men, but can imitate the voices
of cattle in addition. Like Aristotle, Pliny claims that the manticore
does not posses gums in either of its jaws and that its teeth are one
continuous piece of bone.

Ctesias' description includes that the Manticore had a triple row of
teeth which fit into each other like a comb, the face and ears of a
human, blood colored eyes, the body of a lion and a stinging tail
similar to a scorpion. While Pliny the elder argued that it imitated the
voices of men, Ctesias tells us that its voice resembles a union of a
flute and trumpet, that it is exceedingly swift, and has an appetite for
human flesh {[}p.~280{]}.

Aristotle questions Ctesias' account, but nevertheless preserves it in
his own writings. He also reproduces Ctesias' account, that the
Manticore possesses a double rather than triple row of teeth, has
similar feet as a lion and is similar in size to a lion. What is perhaps
important are the anthropomorphic descriptions including those of the
ears and face of a human, and that it has non-anthropomorphic ones such
as grey eyes and red body and possesses a tail like that of a scorpion.
{[}p.~30-31; also see Aristotle, History of Animals 8.28,
605b22-606a10{]}

When we talk about aesthetic function, it is often not an animal's
status as real or imaginary that is a central feature of the animal's
inclusion. Rather, the animal's inclusion into the beastiary is an
aesthetic function. However, it would not be entirely accurate to say
that this function did not possess any truth values. On the one hand, it
may be tempting to say that the value of the animal's inclusion is a
symbolic gesture. But often times the inclusion had Christological
purposes, such as the representation of some important religious
derivation.

There is a stark division in the kinds of reflections these involved as
some animals represented Christian reflections of salvation whereas
others represented threats to salvation. One such example includes
images of a lion breathing on his dead cub and after three days revive
them just as Christ died on the cross and rose after three days
(Morrison and Grollemond 2019, 4). Another point of stark contrast
included animals as the representation of either virtues or vices that
are thought to be common to human experience. A dragon who could kill
elephants by smoothering them with its tail is one such example.

There are several ways in which some art objects may represent more
complex realities. One may think that the purpose of a painting is to
reflect the experiences one had during a war, or peace time event. Such
a painting may express the given facts of the case, for instance an
artist could paint the signing of the Magna Carta.

If the painting were painted by an artist who was alive to witness the
signing of the Magna Carta, then the artist would merely be representing
contemporary reality. If the artist were not alive at the time, then we
could argue that the artist is representing historical reality. The
artist is relying on trusted and well known sources would be different
from the artist imagining an event in which not many direct sources
exist. But an artist can represent a projection of reality as well when
the reality they want to convey is obscure, not well known, highly
contested or hypothetical.

Moral (2021) argues that art's purpose in the middle ages is to
represent reality. It is one thing to accurately and factually represent
a tree, however, Pinero is talking about more complex reality,
especially moral reality. One example of medieval representations of
moral reality include the Persian Manticore. The manticore as depicted
in medieval art is said to represent evil.

According the this definition of medieval art, art implies metaphysical
aspirations and normative elements (Also see De Bruyne 1959). Its
commitment to reality is undeniable, yet the reality which it attempts
to represent is so complex and vast that the limits between what is and
what is not, between what we believe and what we know, between what we
are and what we wish to be, are often blurred (\textbf{cron41?}).

Here, the artist can represent facts as they believe them to be, or as
they wish them to be. It is this latter motivation that Pinero suggests
motivated medieval artwork representing the manticore. There is a
difference here then between modern and medieval artistic function. This
in itself is uninteresting. However, what may be more interesting to the
reader, is the further claim that this difference can be applied in
response to the contemporary account of epistemic injustice.

Owing to the fantastical nature of these accounts, it is highly probable
that they are either misguided representations of some more familiar
animal, or motivated by purposes other than truth. In the classical era,
it is most likely the case that the Manticore is a confused
representation of some other more familiar animal, while in the medieval
period, the motivation for the representation lies outside of truth.

We might think that medieval art has no place in a discussion about
epistemic injustice, but the depiction of the fictional beast called the
manticore, as a personification of deception lends credibility to the
claim that medieval art is an early representation of the dangers of
epistemic harms. According to the modern function of art summed up by
Gilmore (2020), art's function is oriented to ends such as pleasure,
entertainment, and absorption {[}132{]}. Conversely, the medieval
function of art is directed at more orthodox ends intending to motivate
positive emotion states, or desirable behaviors and behavioral
responses.

According to some hermeneutic views centering on epistemic injustice,
because of a misappropriated power imbalance, a person holds false
beliefs about themselves. These beliefs stem from unjust environments
where the individual lacks access to resources that would lead to
accurate self-understanding and moral beliefs. Such beliefs are rooted
in irrational actions performed by that individual.

One such concept of hermeneutic justice has been popularized by Miranda
Fricker (2007) though it was also developed by earlier feminists. As an
artistic object, the Manticore is a fictionalized beast that was born in
Persian art and literature sometime around its cultural zenith. It at
once symbolizes evil, death and hell, but in medieval religious art, it
had come to symbolize the power of Christ over death and hell.

\chapter{Currie's Fiction}\label{curries-fiction}

Why do we think that fairy tales are fiction but news stories are not?
Is fake news the same thing as a fairy tale? During the pandemic, fake
news became a widely discussed talking point. A less widely discussed
talking point however was why don't we just call it fiction? A
derivative of this question is how fiction is made possible. Meaning,
what are the creative differences between fake news, news news and fairy
tales? To answer this question, Gregory Currie (1990) suggests three
possible answers.

According to the front page news, the terrorist leader Amelia Boyd was
killed by the Headhunters Gang. First are the characteristics that
distinguish fiction from non-fiction. This question is a close
investigation of the entities that are in fiction and make them
possible. For instance, the psychological and linguistic resources that
producers and consumers of fiction must have access to.

First, consider that fiction uses language. But there is no universal
linguistic feature shared by all fictional works and there is no
available proof/s confirming what is a fiction. We usually think that
reading a work is an act of reading words and sentences that make it
up.2 However, consider a mime? Or silent films? These are works that do
not use language in the traditional sense, yet we consider them
fictional. Therefore, what makes something a fiction cannot necessarily
be a linguistic feature. Therefore what makes something a fiction is
neither a necessary nor sufficient condition for being fictional.
Finally then, something can be a fiction because of its verbal
structure, but what makes it a fiction does not lie in the text itself.

\section{Semantic Properties}\label{semantic-properties}

Currie concludes that fictionality is a relational property, it is
something possessed in virtue of the text's relations to other things.
Therefore while it is possible that fictionality is a linguistic matter,
what makes it a fiction lies in some other property. (a method of
communication or other communicative practice)

On this account, fictional properties, like non-fictional ones, have the
semantic properties of truth value and reference. What this means is
that a fictional statement can be true and false. For instance, consider
the statement ``Killing an NPC in monitored space will earn you a crime
state''. According to Currie, this statement is true. What makes
something true or false in a fiction, is its relation to the outside
world and so what is important then, is what this relationship is.

Because fiction is a relational property, they are community based and
``depend on the prevailing attitudes of the community, or of a subgroup
of the community rather than the ideas of a single individual'' (Currie
1990, 10).

\begin{quote}
For I take it that an act of joint authorship is exactly that: an act
engaged in by more than one person rather than several distinct acts
undertaken individually and patched together. (Currie 1990, 11)
\end{quote}

The process of creating a fiction begins with the intention. It is not a
singular intention however, but is a social enterprise. So maybe the
reason that we do not call fake news, fiction, is because of the
purposes and intent of fake news. It is intended to deceive rather than
express true propositions. But if this is true, then why do we not
consider that fiction is intended to deceive? Why does fiction have
semantic properties such as truth whereas fake news does not? Currie
(1990)

Owing to the fantastical nature of these accounts, it is highly probable
that they are either misguided or motivated by some purposes other than
truth. In the 57th question of the 1st of the second part of the Summa
Theologica, Thomas Aquinas defends the view of art as a form of right
reason. There is a distinction here between the right reason for making
objects, and the right reason for doing good actions. We should first
focus on the point that is important to both of these activities, is the
right use of one's cognitive faculties. The difference lies in the
produced object, whether an action or thing.

If we are to ask ourselves what justifies our emotional responses to
some given experience, on the one hand we might consider that the
validity of the experience rests in something like the authenticity of
the experience. Assuming that artistic experiences are driven by the
same norms as other experiences, the image of the Manticore would
generate fear in those who were lucky enough to witness an artistic
representation of one, whether literary or visual. However, artistic
experiences are not always governed by the same norms as other kinds of
experiences. Some additional norms that govern our experiences with
fiction include pleasure, entertainment, and absorption (Gilmore 2020,
132). However, this would mark an important divergence from the
functions of art accepted by the medievalist, and those accepted by us
today.

Considering a product that is produced, the desired outcome of the
activity of art, it is not clear what it is exactly, that we care about
in the case of art. When it comes to what is produced in the case of
art, it is not clear whether the final thing produced is an action, or a
material commodity. The common contemporary understanding considers
various perspectives. However, it is important to understand two
important points here.

Moral (2021) argues that this is because art's purpose in the middle
ages is to represent reality. It is one thing to accurately and
factually represent a tree, however, Pinero is talking about more
complex reality, especially moral reality. One example of medieval
representations of moral reality include the Persian Manticore. The
manticore as depicted in medieval art is said to represent evil.

There are several ways in which some art objects may represent more
complex realities. One may think that the purpose of a painting is to
reflect the experiences one had during a war, or peace time event. Such
a painting may express the given facts of the case, for instance an
artist could paint the signing of the Magna Carta. If the painting were
painted by an artist who was alive to witness the signing of the Magna
Carta, then the artist would merely be representing contemporary
reality. If the art were not alive at the time, then we could argue that
the artist is representing historical reality. If the artist is relying
on trusted sources that are well-known, then this would be different
than if the artist were imagining an event in which not many direct
sources exist. But an artist can represent a projection of reality as
well when the reality they want to convey is obscure, not well known,
highly contested etc.

Why do we think that fairy tales are fiction but news stories are not?
Is fake news the same thing as a fairy tale? During the pandemic, fake
news became a widely discussed talking point. A less widely discussed
talking point however was why don't we just call it fiction? A
derivative of this question is how fiction is made possible. Meaning,
what are the creative differences between fake news, news news and fairy
tales? To answer this question, Gregory Currie suggests three possible
answers.

According to the front page news, the terrorist leader Amelia Boyd was
killed by the Headhunters Gang. First are the characteristics that
distinguish fiction from non-fiction. This question is a close
investigation of the entities that are in fiction and make them
possible. For instance, the psychological and linguistic resources that
producers and consumers of fiction must have access to.

First, consider that fiction uses language. But there is no universal
linguistic feature shared by all fictional works and there is no
available proof/s confirming what is a fiction.1 We usually think that
reading a work is an act of reading words and sentences that make it
up.2 However, consider a mime? Or silent films? These are works that do
not use language in the traditional sense, yet we consider them
fictional. Therefore, what makes something a fiction cannot necessarily
be a linguistic feature. Therefore what makes something a fiction is
neither a necessary nor sufficient condition for being fictional.
Finally then, something can be a fiction because of its verbal
structure, but what makes it a fiction does not lie in the text itself.

Currie concludes that fictionality is a relational property, it is
something possessed in virtue of the text's relations to other things.
Therefore while it is possible that fictionality is a linguistic matter,
what makes it a fiction lies in some other property. (a method of
communication, communicative practice)

Fictional properties, like non-fictional ones, have the semantic
properties of truth value and reference. What this means is that a
fictional statement can be true and false. For instance, consider the
statement ``Killing an NPC in monitored space will earn you a crime
state''. According to Currie, this statement is true. What makes
something true or false in a fiction, is its relation to the outside
world. What is important however, is what this relationship is.

The difference lies in the produced object, whether an action or thing.
But here, we should pause momentarily or we risk conflating two senses
of right action. In the ancient context, right action would more
probably be used to reference actions that align with a dominant
organizational structure, for instance the Catholic church. In our own
modern context, we are less apt to recognize a singular dominant moral
authority. This presents a problem since we would want to say something
more akin to `beneficial', or `productive' action where `benefit' and
`product' are more expressive of individual tastes and desires.

Considering a product that is produced, the desired outcome of the
activity of art, it is not clear what it is exactly, that we
collectively care about in the case of art then. When it comes to what
is produced by artistic endeavor, it is not clear whether a normative
account should include a material commodity, or a desired outcome apart
from a material commodity. When we think about what a good artistic
endeavor entails, in the modern conception, there is typically some
monetary component. For instance, an admirable artistic endeavor at the
center of an advertising campaign would ideally result in increased
revenue for the advertising clients and increased business for the
advertiser.

However, the medieval conception of an admirable artistic endeavor
consists of outcomes which invoke some kind of behavior. Whether
dispositions or actions, ultimately there is a desire for an intended
outcome as it correlates specifically with a human actor. Harrison
(2011) argues for instance, that Elizabeth Gaskell's novel Mary Barton:
A Tale of Manchester Life was in part responsible for the UK's expanded
social safety net. Gaskell represented the suffering and distress
experienced by the working poor in a way that allowed her upper-class
Victorian audience to empathize with the lower social classes (Harrison
2011, 258). For instance, we might think that a good artist produces
items of great monetary value. But using the medieval depiction as
expressed above, it seems that a good artist instead motivates an actor
to behave in a way that produces good outcomes. The common contemporary
understanding considers various perspectives. However, it is important
to understand the two important points here.

It is this distinctive kind of epistemic harm/injustice that consists of
undermining our ability to convey knowledge to others and make sense of
our own experiences (Fricker 2007).

Susan Pepper-Gates in a forth-coming paper writes a very engaging and
important paper article applying Miranda Fricker (2007)'s analysis of
epistemic injustice to contemporary Western Religion WR, and the
Philosophy of Religion, PoWR. She argues that the philosophy of religion
should cultivate epistemic-ethical virtues in order to recognize and
include philosopers from minority voices. Doing so helps WR and PoWR
understand and recognize the ways that power shapes our lives as
embodied social beings.

I think that the initial conclusion is a strong one. Namely the ways
that power shapes our lives as embodied social beings is important for
concepts that we associate with who we are. However, I think that the
most important conclusion here is a premise conclusion statement. Rather
than WR and PoWR cultivating epistemic virtues, WR and PoWR should focus
on the cultivation of epistemic virtues, meaning that they should focus
on how epistemic virtue is created. This will give the project of WR and
PoWR two convergent aims.

\begin{enumerate}
\def\labelenumi{\arabic{enumi}.}
\tightlist
\item
  First, it would be to draw attention to the ways in which power
  \textbf{\emph{can}} shape our lives as embodied social beings.
\item
  Cultivating virtue in this context then, means that WR and PoWR ought
  be formative, cultivating virtue in their audiences.
\item
  However, possessing virtue in this sense, makes it less likely that
  power will shape the lives of WR and PoWR's audiences.
\end{enumerate}

\section{Some Questions}\label{some-questions}

A central question here is one of difference in kinds of harm.
Specifically what kind of harm is being done. Additionally, is the kind
of harm that we care about adjudicated solely in virtue of racial, or
ethnic distinctiveness?

\section{Peppers-Gates Analysis of
Fricker}\label{peppers-gates-analysis-of-fricker}

Peppers-Gates first points out that Western Religion, WR, and The
Philosophy of Western Religion, PoWR, have historically ignored work by
intersectional feminist philosophers and philosophers of race who
challenge the traditional approach. However, how is this done? First,
traditional epistmology, TE, focuses on abstract universal knowledge
while the philosophy of religion has traditionally grounded itself in TE
by focusing on rational belief.

In order for the philosophy of religion to cultivate epistemic-ethical
virtues, it should recognize and include philosopers from minority
voices. Peppers-Gates argument takes the following form:

\begin{enumerate}
\def\labelenumi{\arabic{enumi}.}
\tightlist
\item
  Traditional epistmology focuses on abstract universal knowledge.
\item
  Philosophy of religion grounds itself in traditional epistemology and
  focuses on rational belief.
\item
  Therefore WR and PoWR embrace narrow methodologies that exclude others
  and contributes to the marginalization of minority voices.
\item
  Given its focus on abstract universal knowledge and its
  marginalization of minority voices, WR and PoWR ignore work by
  intersectional feminist philosophers and philosophers of race who
  challenge the traditional approach and argue that power shapes our
  lives as embodied social beings.
\item
  As such, WR and PoWR conflate the truth about the existence and
  attributes of a {[}pateriachal{]} God with traditional epistemic
  approaches grounded in {[}narrow methological approaches to
  knowledge{]}.
\item
  By focusing on narrow methodologies and abstract conceptions of
  knowledge, WR and PoWR fail to consider ways that power shapes our
  lives as embodied social beings.
\end{enumerate}

A central point in the above argument is that WR and PoWR have silenced
marginalized groups. These include women, people of color, disabled
persons and members of the LGBTQI+ community. As stated above, the
reason for this is because traditional epistemology focuses on abstract
conceptions of knowledge and thereby fails to consider the ways that
power shapes our lives as embodied social beings. We might think
however, that the argument rests on a notion of rules bias. Namely that
by using certain conversational rules in philosophical discourse, that
PoWR and WR silence minority groups due to their falling in some way to
abide by these rules of discourse.

To see how this is the case, we can use the grammar contrast. Here, we
assume that a person can learn grammar and we exclude them from the
larger conversation until they do. But this would be an unfair response
to the above argument. On the one hand, we could argue that it is
grounded in a misreading of the argument, that minority voices are
excluded from the philosophical lexicon not because they have failed to
abide by the rules of philosophical discourse, but because either 1)
their topics of concern fall outside the margin of interest, or 2)
because of some varying identifying physical characteristic, such as
failure to develop a `rational' argument as determined by the implied
narrow rules of discourse. But this would be a misrepresentation of the
above argument.

The way I understand Peppers-Gates argument, is that marginalized groups
have been excluded from philosophical discourse because of certain
hermeneutic challenges, namely the way they are \emph{understood} as
having failed to abide by the rules of philosophical discourse.
Therefore, we can amend the above argument in the following way:

\begin{enumerate}
\def\labelenumi{\arabic{enumi}.}
\tightlist
\item
  TE focuses on abstract conceptions of knowledge.
\item
  PoWR grounds itself in TE by focusing on `rational' belief in its
  debates about God's existence and attributes.
\item
  Therefore PoWR ignores feminist philosophers and philosophers of race
  who challenge PoWR's rules which are grounded in TE.
\item
  As such, WR and PoWR have silenced marginal groups.
\item
  PoWR's embrace of its narrow methodology and range of topics then is
  based in an exclusion of `others' from meaning making in institutional
  philosophy.
\end{enumerate}

The final premise then says that the narrow methodology employed by PoWR
is not a natural methodology, but is rather exclusive in its
construction. Meaning that it is narrow specifically for the purposes of
excluding minority views and opinions within the lexicon. As such, I
understand that the paper defends the following conclusion:

\begin{quote}
The Philosophy of Western Religion focuses on rational belief and
conflates this with evolving and misapplied lexical rules in its debates
about the existence and attributes of a patriarchal God.
\end{quote}

\section{Silencing}\label{silencing}

Western religion and the philosophy of religion have silenced
marginalized groups by focusing on abstract conceptions of knowledge
which fail to consider ways that power shapes our lives as embodied
social beings. This is because traditional epistemology focuses on
abstract conceptions of knowledge. These abstract conceptions of
knowledge are inadequate because they do not analyze the ways power
shapes our lives as ``embodied social beings''. We are born in webs of
meaning and social relations of power and privilege. Some voices are
heard, meaning they pass review, and therefore shape the social
imagination. Others are ignored however. These others then are incapable
of shaping the social imagination.

On the one hand, a defense of the traditional power structure would make
sure to emphasize what it means to participate in the sort of debate
that does have the potential to shape the social imagination. How is
one's voice heard? One way is through the enforcement of so-called
grammar rules. Taking an art class, the professor attempts to
communicate the rules which apply to a given art genre. However, the
students note the many deviations from these rules, the supposed
exceptions. The professor then replies by saying that ``once you know
how to apply the rules, then you are allowed to break the rules''. An
imagined analogy would be a reviewer rejecting a paper for the alleged
abuse of grammar rules that a more seasoned author would be allowed to
bend. A concern here is one where the so-called rules are being
inappropriately applied as a means of gate-keeping who gets to shape the
`social-imagination'. A student who fails to learn the rules, is not
accepted into the art community, where work is posted in galleries, sold
at auctions, is not publicly (positively or negatively) reviewed etc,.
Of additional importance here, is that assuming a more seasoned author
is justified in bending the lexical rules of philosophical discourse,
than the more seasoned author can also be seen as contributing to the
development of the lexicon rules surrounding philosophical discourse. As
such, the more seasoned professional is allowed to shape the rules of
discourse, and thereby contribute to the social imagination.

As such, the harm which ought to concern us here includes the exclusion
of individuals from the generally shared social understanding. Excluding
voices in this way contributes to their harm or can be used as a weapon
to harm individual. But what is this weapon and how specifically does it
as a weapon cause harm? In the following section, I will show how
participation in the shared lexical structure has been called a marker
of humanity. Therefore gatekeepers undermine the process by which a
person becomes human.

\section{Hermeneutic Injustice}\label{hermeneutic-injustice}

\begin{enumerate}
\def\labelenumi{\arabic{enumi}.}
\tightlist
\item
  When DSS (Dominant Social Structure) of meaning exclude you or distort
  your experiences, and no social concepts that you recognize capture
  your vision of yourself, then your capacity as a knower and your own
  self-understanding have been stunted.

  \begin{enumerate}
  \def\labelenumii{\arabic{enumii}.}
  \tightlist
  \item
    But this introduces a very narrow vision of social structure. Isn't
    it possible that while you are excluded from a dss, a non-dss, one
    that includes you and does not distort your experiences, could still
    capture a vision of yourself, and your capacity as a knower and
    fortifying your self-understanding.
  \item
    What is needed then is a new non-dss, education.
  \item
    While a person is excluded from a dss, they could have their
    capacity as a knower fortified by a non-dss, education reimagined.
  \end{enumerate}
\end{enumerate}

The need to reimagine eduction.

\begin{enumerate}
\def\labelenumi{\arabic{enumi}.}
\tightlist
\item
  WR and PoR become DSS's of meaning and exclude some or distort their
  experiences
\item
  But this introduces a very narrow vision of social structure.

  \begin{enumerate}
  \def\labelenumii{\arabic{enumii}.}
  \tightlist
  \item
    It is possible that while you are excluded from a dss, a non-dss,
    one that includes you and does not distort your experiences, could
    still capture a vision of yourself, and your capacity as a knower
    and fortifying your self-understanding.
  \item
    So while a person is excluded from a dss, they could have their
    capacity as a knower fortified by a non-dss.
  \end{enumerate}
\item
  Therefore 2 this does not mean that there are no social concepts that
  you recognize and capture your vision of yourself and your capacity as
  a knower and your own self-understanding are stunted but not by the
  DSS, but the lack of the non-dss.
\end{enumerate}

\subsection{One Possible example of hermeneutic injustice against
religious
minorities}\label{one-possible-example-of-hermeneutic-injustice-against-religious-minorities}

Secular community marginalization against people of religious faith.

\begin{enumerate}
\def\labelenumi{\arabic{enumi}.}
\tightlist
\item
  PoRF may lack 'sensibilities and resources needed to make sense of
  religious people's experiences,'' (Kidd, 20217).
\item
  But does not capture gravity of systemic cases of hermeneutical
  injsutice whose targests experience multiple forms of injustice.

  \begin{enumerate}
  \def\labelenumii{\arabic{enumii}.}
  \tightlist
  \item
    Hermeneutical marginalization ``entails mariginalization of a
    socioeconomic sort', targets will not 'participate on equal terms
    within the practices of law, academia, politics, religion,
    journalism, media, etc. (Fricker, 2010).
  \item
    However, DSS vs.~non-DSS? While DSS of meaning may distort
    experiences, non-dss may not, instead rejecting dss of meaning which
    says that ``there are no social concepts which capture a religious
    person's vision of themselves (i.e., made in image of God, human
    dignity bestowed by God, not DSS, etc.)''. Consider the following
    claims:
  \item
    Here, Fricker, a published philosopher, the DSS, says persons of a
    certain sort are marginalized and therefore have been silenced. But
    here, marginalization means they lack a marker of humanity. As such,
    Fricker claims that certain individuals lack a marker of humanity.
  \item
    Why? For not adhering to rules of rationality practiced by the dss.
    Surely they could however do so?
  \item
    Because of 1, they lack participation in shared meaning.
  \item
    Therefore lack capacity as knowers.
  \item
    From 3, their self-understanding is mitigated and stunted.
  \item
    Participation in shared hermeneutic, i.e., 1-4, in shared conceptual
    structure is a marker of humanity'',
  \item
    Therefore, because of exclusion from shared conceptual structure,
    they lack a marker of humanity.
  \end{enumerate}
\end{enumerate}

but rather than `victims' being excluded from shared conceptual
structure and lacking capacity as knowers and lacking marker of
humanity, they instead are included in different conceptual structure
(e.g., religious conceptual structure) which fortifies their capacity as
knowers and self-understanding, acting as marker of humanity.

\chapter{On Greek Tragedy}\label{on-greek-tragedy}

\section{Three Kinds of Tragic
Conflict}\label{three-kinds-of-tragic-conflict}

\begin{enumerate}
\def\labelenumi{\arabic{enumi}.}
\tightlist
\item
  Good people being ruined because of things that just happen to them,
  things they do not control.
\item
  Good people doing bad things, things opposite of their ethical
  character and commitments because of things they do not control.
\item
  Good people doing bad things without direct physical compulsion and
  opposite of their ethical character and commitments though done in
  light of circumstances of equally bad choices.
\end{enumerate}

Agammenon is waiting on the shore with his entire fleet. He intends to
carry out a directive from Zeus to avenge his brother whose wife ran
away with Paris, prince of Troy. However, because of an earlier
transgression against Artimis, Artimis is holding back the wind,
stranding his fleet. Agammenon, as the result of an omen of an eagle
devouring a hare, is told that in order to appease Artimis, he must
sacrifice his daughter Iphigenia.

As such, Agammenon is faced with two impossible choices, fail to carry
out Zeus's command, or kill his daughter. It is important to point out
that the fleet is already facing starvation being stranded on the coast.
However, Agammenon's sin, is not that he choices sacrificing his
daughter, instead it is the way he sacrificed his daughter.

\begin{itemize}
\tightlist
\item
  Like in the biblical story of Abraham, a good man must kill an
  innocent child ``out of obedience to a divine command'' {[}35{]}.
\item
  Agammenon's sin however, is not the choice, but rather how the choice
  was made:

  \begin{itemize}
  \tightlist
  \item
    Agammenon now begins to cooperate inwardly with necessity, arragning
    his feelings to accord with his fortune.
  \item
    From the moment he makes his decision, itself the best he could have
    made, he strangely turns himself into a collaborator, a willing
    victim {[}35{]}.
  \end{itemize}
\end{itemize}

Progression of loss of agency:

\begin{itemize}
\tightlist
\item
  Necessity of the act
\item
  to the rightness of the act
\item
  to the rightness of the feelings of the act
\end{itemize}

In another paper, I highlighted some concerns raised by Wendy Laybourn
(2018) regarding representation and authenticity. Here, I want to draw
attention to some additional similar concerns. What is important here,
is that these criticisms unfairly focus too narrowly on the black
experience and do not do a good enough job in recognizing that these
experiences are universal. Given their universality, I argue that a best
response is found in a Christian theodicy to evil as a form of
suffering.

Laybourn (2018)'s a criticism of black authenticity and what she calls
commodified realness is a commodified identity that is sold to a broader
public as black authenticity. This criticism is centered around the
effects of what she calls controlling images in popular media. Of deep
concern here, is that many of the members of the communities depicted by
such images consider aspects of these images to be authentic
representations and not in themselves harmful. Yet too, this is also a
common trope.

I present this case as an example of a paradigmatic concept of
hermeneutic injustice. It is not clear just what aspects of the member's
lives are represented, nor specifically which images are harmful if such
images are accepted as representative. An important question here has to
do with the relationship between popular narratives that contain these
images, and personal identity. One approach argues that we are defined
by the stories we accept about ourselves, but this is controversial and
it is the point I will address.

\subsection{Wendy Laybourn's Cost of Being
Real}\label{wendy-laybourns-cost-of-being-real}

Wendy Laybourn (2018) draws a comparison between black authenticity and
what she calls commodified realness. According to Laybourn, identity is
sold to a broader public as black authenticity. This criticism is
centered around the effects of what she calls \emph{controlling images}
in popular media. Of deep concern here, is that many of the members of
the communities depicted by such images, consider aspects of these
images to be authentic representations of themselves and the communities
in which they live. Therefore, there is a public perception that such
images are not in themselves harmful. If they are in some degree
controlling, then they would be harmful, but there is strong
disagreement regarding whether or not they are (20C History Project
2014; Collins 2006, 2009; Sartwell 1998).

\chapter{Patricia Collins Black Feminist
Thought}\label{patricia-collins-black-feminist-thought}

\phantomsection\label{collins-knowledge-power}
In the 2009 version of Collin's Black Feminist Thought, BFT, Collins
draws greater emphasis between relationships of knowledge and power. It
is often said

\begin{itemize}
\tightlist
\item
  that knowledge = power,
\end{itemize}

but a worrying theme throughout the book, is that this is wrong and that
rather

\begin{itemize}
\tightlist
\item
  power = knowledge.
\end{itemize}

One might notice the misapplication of a logical principle here, because
whenever we use the equals sign, there is not a logical difference
between power = knowledge and knowledge = power.

\begin{quote}
Rather, this says that power is identical in value to knowledge and if
power is identical to knowledge, then knowledge is identical to power.
\end{quote}

In other words, the initial application claims that the direction of fit
is not directional but rather than uni- or bi-directional.

For instance, we could look at some possible relationships between
knowledge and power below:

\begin{itemize}
\tightlist
\item
  knowledge \(\rightarrow\) power
\end{itemize}

or

\begin{itemize}
\tightlist
\item
  power \(\rightarrow\) knowledge
\end{itemize}

or

\begin{itemize}
\tightlist
\item
  knowledge \(\leftrightarrow\) power
\end{itemize}

Initially, we might think that assuming that if one has knowledge then
the correct state of affairs should be that they will fulfill their
aims, whatever those may be. This is often a selling point for higher
education.

\phantomsection\label{collins-struggle}
This is a common position in what is lauded as ``African American
Studies''. Often times these positions are wrapped in the language of
``struggle''. While Collins notes that this ``struggle'' is a struggle
for empowerment, it is not particularly clear who the struggle is
directed at. She argues that the struggle is global and transcendent,
and she defines it as one for ``survival.'' But she expands the target
in the most recent edition of her book to include recent developments in
the areas of sexuality and ``sexual politics'' of black womanhood.

Some similar examples of this criticism can be found in Zakiyyah Iman
Jackson (2020) or Mbembe (2017)'s Critque of Black Reason. Jackson
(2020) defends the view of race as a social construct. Here black
represents a void of reason and a surplus of animality. It is not clear
what the counter to this claim should be as it is not represented
anywhere in either work. However, looking at what is deemed black
culture broadly in popular media, this counter claim is even more
difficult to understand given that many African Americans fail to
explicitly and publicly criticize, or at the very least, argue that such
behaviors are not culturally representative.

Popular artists who reflect similar characteristics as those deemed
synonymous to the black race are praised and vaunted. Black female
artists for instance, who reinforce stereotypes of blacks as
hypersexualized, vulgar, using codes such as twerking to exemplify their
over hyper-sexualized nature, are praised as empowering or authentic,
where as black male artists are praised as being ``real'' or ``gritty''
for expressing unjustified violence directed at other blacks. But this
is just to mimic the age old racist justifications for slavery although
slavery, in name, has been outlawed. According to this justification,
blackness represents a blurred distinction between human and animal.

Man distinguishes himself from animality, but this is not the case for
the Black Man, who maintains within himself, albeit with a certain
degree of ambiguity, animal possibility. A foreign body in our world, he
is inhabited­under cover---by the animal. To debate Black reason is
therefore to return to the collection of debates regarding the rules of
how to define the Black Man: how he is recognized, how one identifies
the animal spirit that possesses him, ­under which conditions the ratio
penetrates and governs the animalitas. (Mbembe 2017, 56)

It is not necessarily clear where to find the literature on specifically
what is at issue here. Because of this, I think it will be helpful to
start with something of an analogy, specifically one from the popular
narrative. In the newest Wicked screen adaptation, there is an
underlying story line. The protagonist, the `wicked' witch finally comes
to accept that she will never be accepted by the broader public and as
such, will always stand out as a scapegoat of sorts owing to her dark
green skin. Therefore she fully leans into her `wicked' persona.

However, perhaps a detail which could not be included as prominently in
Gregory Maguires 1995 novel. The amazing actress, Cynthia Erivo, is an
African American with especially prominent features, accentuated owing
to the excellent applied makeup and superb acting by Evrivo. Through
both movies, it is obvious that she isn't green in reality, but African
American. We see the story of the witch as a case of othering. As such,
especially poignant was the moment when Dorothy, believing the witch can
be killed throws a bucket of water on her, an old racist trope about
African Americans and bathing.

It is a fortunate accident I suppose, director John M. Chu taking
perfect advantage of the fact that the witch in the original dies having
had a bucket of water thrown onto her. The witch, fully adopting her
role, seemingly for the good of the people of Oz, `dies', but in the
fiction, merely goes into hiding. Ultimately Elphaba and Prince Fiyero,
turned Straw man, strike out to the desert after her sacrificial death
accepting that they will forever be scapegoats should they choose to
stay in Oz.

This is the naive view:

It can perhaps be with a bit of irony that one pursues the study of
African American literary critiques. At once, these critiques are
leveled against a system that is responsible for their own publishing.
Frantz Fanon's, we might naively say, is one such example. It is perhaps
uncouth to look at a book's publishers in an effort to more fully
understand the semantic content of that book. Explanations of how the
book came to exist are not thought to lend much support in assessing the
validity of an author's claims. Unfortunately however, the temptation to
do so can be quite overwhelming at times.

Perhaps there are reasons why an understanding of the processes
surrounding a book's existence are an important part of the existence of
that book, for instance the Western education of the author, the book's
distribution, the retail outlets, libraries, and classrooms, that
represent the demand for the book.

Naively, I might be led to think that without any of these artifacts of
a culture, that a given book would not exist. More important however,
are those specific motivations that go into the book's creation. The
demands of the editor, positive reviews by reviewers, a public's
insistence that such a book needs to be written. Lastly, there is a
worry that only certain authors can write a given book, such that if
anyone else where to write it, it would not be received by publishers,
editors, or the public.

A similar critique here is Frantz Fanon's exploration of the effects of
colonialism in the psychology of Africans. Of concern both for Fanon's
work as well as Jackson (2020), it can be difficult to disentangle what
might be considered a person's culture in isolation from the affects of
colonialism. I myself am sceptical that such a project could be
successful.

To read a work like Fanon's \emph{Black Skin, White Masks}, then, for
someone as unsophisticated as myself, can be quite a daunting task. All
I see are continuous contradictions between the `criticisms' of `western
imperialism' and the words themselves which to my unsophisticated ears
sounds as the very words used by those `western imperialists', at least
the ones they themselves would use. In the foreword, written by Anthony
Appiah, it is claimed \emph{Masks} is confronting the psychological
affects of racism.

This is an important point, given that Fanon himself is trained as a
psychologist. However, Fanon begins the introduction using colonial
expressions of old world notions of blackness as objects devoid of
personhood. As such, he himself is doing depriving others of personhood
and I am not sure why he is doing so. These expressions prove to me to
be quite difficult obstacles to surmount. The `educated elite' who will
undoubtedly be well-versed in this literature by the time my own
children begin to sit in their classrooms to my unsophisticated mind
will have their own heads filled with the most awful views and ideas
about who my children are. How are they supposed to be able to draw from
their bowels the ability to cognitively navigate around their tutors
words that undoubtedly will have been shaped by these horrendous views?

The books target initially then is ``black'' empowerment which would
involve the acquisition of knowledge. The problem with the above
however, is that any correlation between knowledge and power is wrong
from a Christian world view. Rather one is grounded in power if they
have humility. ``The fear of the Lord is the beginning of wisdom.''
Collins mentions the 19th century black feminist scholar Maria Stewart,
noting that she urged self-reliance and independence and that her work
was suppressed along with many other black feminist scholars of the era.
However, we as a society know a lot more about the importance of
community and inter-personal relationships, a fact which Stewart herself
draws attention to:

This is a common position in what is lauded as ``African American
Studies''. Often times these positions are wrapped in the language of
``struggle''. While Collins notes that this ``struggle'' is a struggle
for empowerment, it is not particularly clear who the struggle is
directed at. She argues that the struggle is global and transcendent,
she defines the struggle as one for ``survival.'' But she expands the
target in this edition to include recent developments of sexuality and
``sexual politics'' of black womanhood.

Here, the implication is not on knowledge, knowledge without wisdom is
useless and a fool's errand, it does not lead to ``empowerment''. As
noted in the feminist and social epistemology literature, there are many
divergent accounts of knowledge and while the discourse may focus on one
account of knowledge as being central to empowerment, they are often
missing other accounts of the importance of knowledge. Collins mentions
the 19th century black feminist scholar Maria Stewart, noting that she
urged self-reliance and independence and that her work was suppressed
along with many other black feminist scholars of the era. However, we as
a society know a lot more about the importance of community and
inter-personal relationships, a fact which Stewart herself draws
attention to:

\begin{quote}
You have souls committed to your charge. . . . It is you that must
create in the minds of your little girls and boys a thirst for
knowledge, the love of virtue, . . . and the cultivation of a pure
heart.
\end{quote}

Here, an important development is that Stewart draws a deep connection
between knowledge and virtue.

\chapter{The Virtue of Luck (a reading of luck and ethics from Fragility
of
Goodness)}\label{the-virtue-of-luck-a-reading-of-luck-and-ethics-from-fragility-of-goodness}

What is the relationship between human virtue and environmental
variables? In order for a trait to be the virtue of a human being, some
would like to say that it has to represent choices, rationality,
actions, beliefs. However, often times these depend on environmental
variables which we do not control. Most importantly among these are the
values that we hold dear to us. Whether or not we recognize the
importance of close relationships for instance. However, if the luck
question is accurate, even these values depend on factors that are
outside of our control. Our fragility, need for loved ones, friendship,
etc., can at one time depend on whether or not we've encountered
sufficient hardships for instance, but even outside the experience of
hardships, friendships, which are also external and an environmental
variable at once hold intrinsic importance, although these are external
to us. As such ``human goodness is nourished, even constituted, by
external happenings''. {[}1-2{]}

\chapter{Non-Scientific Deliberation}\label{non-scientific-deliberation}

\section{Practical Deliberation}\label{practical-deliberation}

\begin{itemize}
\tightlist
\item
  Not scientific
\item
  appropriate criterion of correct choice is a person of practical
  wisdom (thoroughly human being)
\item
  therefore vulnerable to ungoverned \emph{tuche} whereas Plato thinks
  that a good human life, a person exercising practical wisdom is
  invulnerable to luck
\end{itemize}

How then are good judgments of value made?

\subsection{The vulnerability of the Good Human
Life}\label{the-vulnerability-of-the-good-human-life}

All human excellence requires external resources and necessary
conditions.

\begin{itemize}
\tightlist
\item
  External objects to exercise virtues towards
\item
  External objects to receive our virtues
\end{itemize}

The most important of these require relational virtues.

E.g.,

\begin{itemize}
\tightlist
\item
  Good activities of citizenship
\item
  Political attachment
\item
  Personal love
\item
  Friendship
\end{itemize}

\subsubsection{The Relational Harm of
Slavery}\label{the-relational-harm-of-slavery}

Good character requires civic activity, good political surroundings.

A slave however, lacks:

\begin{itemize}
\tightlist
\item
  choice
\item
  \emph{philia}
\end{itemize}

A person needs to be able to make free and un-coerced choices. A slave
is unable to do so since he belongs to someone else who makes choices
for him. Owing to his inability to make free and unconstrained choices,
a slave is unable to use their capacity for practical reason, choosing
objects for their own sake. As such, they cannot share in \emph{philia}
which ``is based on mutual respect for choice and character'' {[}348{]}.

\chapter{Morality as Privilege}\label{morality-as-privilege-1}

Montaque Reynolds (Stetson University)

\hfill\break

\section{Abstract}\label{abstract-1}

This concerns an interesting and popular assertion, that it takes a
certain amount of material wealth to be a good person. I suspect that
this is a project of capitalism: buy our stuff or you will be a bad
person . . .

\chapter{The Virtue of Luck (a reading of luck and ethics from Fragility
of
Goodness)}\label{the-virtue-of-luck-a-reading-of-luck-and-ethics-from-fragility-of-goodness-1}

What is the relationship between human virtue and environmental
variables? In order for a trait to be the virtue of a human being, some
would like to say that it has to represent choices, rationality,
actions, beliefs. However, often times these depend on environmental
variables which we do not control. Most importantly among these are the
values that we hold dear to us. Whether or not we recognize the
importance of close relationships for instance. However, if the luck
question is accurate, even these values depend on factors that are
outside of our control. Our fragility, need for loved ones, friendship,
etc., can at one time depend on whether or not we've encountered
sufficient hardships for instance, but even outside the experience of
hardships, friendships, which are also external and an environmental
variable at once hold intrinsic importance, although these are external
to us. As such ``human goodness is nourished, even constituted, by
external happenings''. {[}1-2{]}

\chapter{On Greek Tragedy}\label{on-greek-tragedy-1}

\section{Three Kinds of Tragic
Conflict}\label{three-kinds-of-tragic-conflict-1}

\begin{enumerate}
\def\labelenumi{\arabic{enumi}.}
\tightlist
\item
  Good people being ruined because of things that just happen to them,
  things they do not control.
\item
  Good people doing bad things, things opposite of their ethical
  character and commitments because of things they do not control.
\item
  Good people doing bad things without direct physical compulsion and
  opposite of their ethical character and commitments though done in
  light of circumstances of equally bad choices.
\end{enumerate}

Agammenon is waiting on the shore with his entire fleet. He intends to
carry out a directive from Zeus to avenge his brother whose wife ran
away with Paris, prince of Troy. However, because of an earlier
transgression against Artimis, Artimis is holding back the wind,
stranding his fleet. Agammenon, as the result of an omen of an eagle
devouring a hare, is told that in order to appease Artimis, he must
sacrifice his daughter Iphigenia.

As such, Agammenon is faced with two impossible choices, fail to carry
out Zeus's command, or kill his daughter. It is important to point out
that the fleet is already facing starvation being stranded on the coast.
However, Agammenon's sin, is not that he choices sacrificing his
daughter, instead it is the way he sacrificed his daughter.

\begin{itemize}
\tightlist
\item
  Like in the biblical story of Abraham, a good man must kill an
  innocent child ``out of obedience to a divine command'' {[}35{]}.
\item
  Agammenon's sin however, is not the choice, but rather how the choice
  was made:

  \begin{itemize}
  \tightlist
  \item
    Agammenon now begins to cooperate inwardly with necessity, arragning
    his feelings to accord with his fortune.
  \item
    From the moment he makes his decision, itself the best he could have
    made, he strangely turns himself into a collaborator, a willing
    victim {[}35{]}.
  \end{itemize}
\end{itemize}

Progression of loss of agency:

\begin{itemize}
\tightlist
\item
  Necessity of the act
\item
  to the rightness of the act
\item
  to the rightness of the feelings of the act
\end{itemize}

\subsection{The vulnerability of the Good Human
Life}\label{the-vulnerability-of-the-good-human-life-1}

All human excellence requires external resources and necessary
conditions.

\begin{itemize}
\tightlist
\item
  External objects to exercise virtues towards
\item
  External objects to receive our virtues
\end{itemize}

The most important of these require relational virtues.

E.g.,

\begin{itemize}
\tightlist
\item
  Good activities of citizenship
\item
  Political attachment
\item
  Personal love
\item
  Friendship
\end{itemize}

\subsubsection{The Relational Harm of
Slavery}\label{the-relational-harm-of-slavery-1}

Good character requires civic activity, good political surroundings.

A slave however, lacks:

\begin{itemize}
\tightlist
\item
  choice
\item
  \emph{philia}
\end{itemize}

A person needs to be able to make free and un-coerced choices. A slave
is unable to do so since he belongs to someone else who makes choices
for him. Owing to his inability to make free and unconstrained choices,
a slave is unable to use their capacity for practical reason, choosing
objects for their own sake. As such, they cannot share in \emph{philia}
which ``is based on mutual respect for choice and character'' {[}348{]}.

\chapter{Narrative Regrets in Truth and
Fiction}\label{narrative-regrets-in-truth-and-fiction}

God as the Best Possible Epistemic Ground of Human Dignity

Following @gold09's short paper \emph{Narrative Thinking, Emotion, and
Planning}, I consider the representation of emotional expressions of
regret, \emph{narrative regret}. I draw specific attention to the
context of the representation of distinct special relationships in
fiction. Important here are discourses on the fictional representation
of such relationships as a special kind of human agency. I argue that a
related concern are the assumptions about what is true and what is
fiction in popular fiction with regard to special relationships, where
what is true are expressions of relational sanctification that are
grounded in a divine union of love appropriate to friendship.

\hfill\break

\section{Testimonial and Hermeneutic
Injustice}\label{testimonial-and-hermeneutic-injustice}

\subsection{Wendy Laybourn's Cost of Being
Real}\label{wendy-laybourns-cost-of-being-real-1}

Wendy Laybourn (2018) draws a comparison between black authenticity and
what she calls commodified realness. According to Laybourn, identity is
sold to a broader public as black authenticity. This criticism is
centered around the effects of what she calls \emph{controlling images}
in popular media. Of deep concern here, is that many of the members of
the communities depicted by such images, consider aspects of these
images to be authentic representations of themselves and the communities
in which they live. Therefore, there is a public perception that such
images are not in themselves harmful. If they are in some degree
controlling, then they would be harmful, but there is strong
disagreement regarding whether or not they are (20C History Project
2014; Collins 2006, 2009; Sartwell 1998).

Research on black authenticity identifies ``oppressive othering'' and
``white standards of beauty'', particularly lighter skin tone, as key
components of commercial success in the entertainment industry.

\begin{itemize}
\tightlist
\item
  Oppressive Othering (commodified black authenticity, cba)
\item
  White Standards of Beauty (wsb)
\end{itemize}

Are rap lyrics a case of oppressive othering? By using rap lyrics,
Laybourn argues that we can compare oppressive othering with white
standards of beauty arguing that each are ways mass media disenfranchies
individuals. However, the way she does so is by drawing attention to a
positive correlation between lighter skin tone and higher Billboard
chart ratings versus linking less dominant social classes with
antisocial behavior.

Laybourn follows Nguyen and Anthony (2014). According to Nguyen \emph{et
al}, the research on black authenticity reveals two trends. The first is
commodification when cultural productions of authenticity affect how
individuals or cultural products fare in the economic market
{[}p.~771{]}. The second includes constructing expectations and
self-identities. Here a similar relation can be seen between rap lyrics
(oppressive othering) and higher chart ratings where higher chart
ratings are predicated when an artists meets particular expectations.
Though at the center of this issue is that these expectations and
identities justify or legitimate membership when they are fulfilled.

Therefore her suggestion seems to be that a harmful representation of
blackness is sold to the public as ``authentic blackness'' by promoting
artists with lighter skin tones. Importantly, this particular commodity
then ``shapes'' what is accepted as authentic blackness, and what is
commodified as authentic blackness. However, the implication here is
that skin tone equates to one's identity.

Therefore, here identity is constructed as authentic when the successful
signification of what is accepted as ``real'' or ``true'' for cultural
products, music etc, and individual's identities (a person as an
authentic artist) is socially constructed. Therefore, Nguyen \emph{et
al} suggest that the research on ``black authenticity'' across contexts
ought to include varying ideals and expectations which affect ``what it
means to be black'' whether in personal, public, and cultural identities
{[}p.~779{]}.

In addition to the ``oppressive othering'' of cba, ``white standards''
of beauty are also sold. But it isn't clear what correlation there is
between cba and wsb if any. Laybourn argues that both wsb and cba stem
from attitudes which position blacks as an outgroup and whites as an
in-group. I think the focus on beauty standards is distracting from what
is truly important about her paper, which are those qualities that we
should consider as the essential nature of a person. In other words, we
should not think that the color of a person's skin is an essential
marker of who they are.

This is the naive view:

It can perhaps be with a bit of irony that one pursues the study of
African American literary critiques. At once, these critiques are
leveled against a system that is responsible for their own publishing.
Frantz Fanon's, we might naively say, is one such example. It is perhaps
uncouth to look at a book's publishers in an effort to more fully
understand the semantic content of that book. Explanations of how the
book came to exist are not thought to lend much support in assessing the
validity of an author's claims. Unfortunately however, the temptation to
do so can be quite overwhelming at times.

Perhaps there are reasons why an understanding of the processes
surrounding a book's existence are an important part of the existence of
that book, for instance the Western education of the author, the book's
distribution, the retail outlets, libraries, and classrooms, that
represent the demand for the book.

Naively, I might be led to think that without any of these artifacts of
a culture, that a given book would not exist. More important however,
are those specific motivations that go into the book's creation. The
demands of the editor, positive reviews by reviewers, a public's
insistence that such a book needs to be written. Lastly, there is a
worry that only certain authors can write a given book, such that if
anyone else where to write it, it would not be received by publishers,
editors, or the public.

A similar critique here is Frantz Fanon's exploration of the effects of
colonialism in the psychology of Africans. Of concern both for Fanon's
work as well as Jackson (2020), it can be difficult to disentangle what
might be considered a person's culture in isolation from the affects of
colonialism. I myself am sceptical that such a project could be
successful.

To read a work like Fanon's \emph{Black Skin, White Masks}, then, for
someone as unsophisticated as myself, can be quite a daunting task. All
I see are continuous contradictions between the `criticisms' of `western
imperialism' and the words themselves which to my unsophisticated ears
sounds as the very words used by those `western imperialists', at least
the ones they themselves would use. In the foreword, written by Anthony
Appiah, it is claimed \emph{Masks} is confronting the psychological
affects of racism.

This is an important point, given that Fanon himself is trained as a
psychologist. However, Fanon begins the introduction using colonial
expressions of old world notions of blackness as objects devoid of
personhood. As such, he himself is doing depriving others of personhood
and I am not sure why he is doing so. These expressions prove to me to
be quite difficult obstacles to surmount. The `educated elite' who will
undoubtedly be well-versed in this literature by the time my own
children begin to sit in their classrooms to my unsophisticated mind
will have their own heads filled with the most awful views and ideas
about who my children are. How are they supposed to be able to draw from
their bowels the ability to cognitively navigate around their tutors
words that undoubtedly will have been shaped by these horrendous views?

As such, cbas are dangerous when they serve as an avenue for interaction
with black culture and a proxy for interpersonal interaction, as
evidenced not only by Fanon, but by an entire generation of scholarship
around what it means to be black. This is true because they undermine
true companionship in that a person is only becoming acquainted with a
false image. This image is especially harmful because it is created
within a framework ``that'' reifies anti-blackness, but it is not only
blacks that are harmed, but rather the human person. In this way, cbas
have no bearing on reality, but they take on many seemingly benign forms
as well. There is a public perception that such images are not in
themselves harmful, for instance as evidenced with the centrality of
Fanon's literature in academia, perhaps this explains the prevalence of
cbas in popular culture, even among those who the images are harming
most. An important consequence of this is that many of the members of
the communities depicted by such images, consider aspects of these
images to be authentic representations of themselves and of the
communities in which they live and of themselves. Therefore, there is
strong disagreement regarding whether or not they are harmful (20C
History Project 2014; Collins 2006, 2009; Sartwell 1998).

In another paper, I highlighted some concerns raised by Wendy Laybourn
(2018) regarding representation and authenticity. Here, I want to draw
attention to some additional similar concerns. What is important here,
is that these criticisms unfairly focus too narrowly on the black
experience and do not do a good enough job in recognizing that these
experiences are universal. Given their universality, I argue that a best
response is found in a Christian theodicy to evil as a form of
suffering.

Laybourn (2018)'s a criticism of black authenticity and what she calls
commodified realness is a commodified identity that is sold to a broader
public as black authenticity. This criticism is centered around the
effects of what she calls controlling images in popular media. Of deep
concern here, is that many of the members of the communities depicted by
such images consider aspects of these images to be authentic
representations and not in themselves harmful. Yet too, this is also a
common trope.

I present this case as an example of a paradigmatic concept of
hermeneutic injustice. It is not clear just what aspects of the member's
lives are represented, nor specifically which images are harmful if such
images are accepted as representative. An important question here has to
do with the relationship between popular narratives that contain these
images, and personal identity. One approach argues that we are defined
by the stories we accept about ourselves, but this is controversial and
it is the point I will address.

Eleonore Stump (2010) writes on the philosophical problem of evil,
arguing that what is truly of concern is not the problem of evil, but
rather the problem of suffering. In order to understand the problem that
evil presents, we must first understand the nature of suffering. It is
in suffering that evil is the most pronounced according to Stump.

Here, an important component of suffering, is psychological dissolution
where one experiences a disentegration of the will. A result of this
disentegration is that a person becomes divided from themselves and is
no longer able to recognize the desires of their own heart. These
desires are ultimately relational. As such, the person who has become
divided suffers in a state Professor Stump refers to as ``willed
loneliness'', which is a state of psychic disentegration.

Unfortunately, I will not be able to give Professor Stump article the
treatment that it deserves. However, before we can even understand what
psychic disintegration is like, it is first necessary to understand why
we might think union with another is a mark of spiritual health. For
this, we need a theory of identity. Towards the end of this paper, I
introduce such a theory. First however, I discuss a number of examples
of what relational desire looks like.

\section{Robust and Generic Expressions of Regret and Other Popular
Forms of
Music}\label{robust-and-generic-expressions-of-regret-and-other-popular-forms-of-music}

The country music singer George Jones was born September 12, 1931 in
Saratoga Texas. He would die at the age of 81 which is considerable
considering the storied life he lived. In the podcast; \emph{Cocaine and
Rhinestones}, and later the book of the same name, Tyler Mahan Coe
(2024) dismantles common tropes associated with the singer. Of special
importance was Jone's popular marriage to the other country singer Tammy
Wynette which Coe accusses of being an industry publicity stunt, and his
marriage to Nancy Pulvado which many credit with Jone's recovery from
drug and alchohol addiction.

The song \emph{He Stopped Loving Her Today} was completed in 1979 by
Bobby Braddock and Curly Putnam, writers for Nashville's Tree Publishing
Company (now Sony Music Publishing) (Jones 1979). At least a year prior
to its completion however, George Jones would be told about it by his
producer Billy Sherrill. It would go on to be iconic of country music
from the era. In it, the narrator speaks of a neighbor's dissolution of
their relationship. It is never made clear what kind of relationship it
is, whether a legally recognized marriage, friendship, love affair etc.
However, what is made manifest throughout the song is its intensity with
the narrator expressing that the day his neighbor died, is the first
expression of joy he'd seen from the neighbor since his lover left.

Love songs may express regret for those moments prior to the narrator's
introduction to the beloved. It is possible that when a song expresses
regret for not having met someone sooner, they are thinking about how
their life might have gone had they married sooner. The marriage itself
is not meaningful for the individual as a stand-alone event. Rather, the
marriage only has causal influence on the subsequent events of that
person's life and it is the person who defines the event in such a way
to make it a meaningful marker of their own identity. It is where the
marriage sits in the overall context of the course of a given
relationship for instance. A marriage event in the context of a short
and harmful relationship has a very different meaning than one in the
context of a long and loving relationship. Most importantly however, it
is important in the context of the meaning that a person derives from
the event.

Many popular songs are saturated in the emotion of regret. In such
contexts, the narrative expresses a negative emotion in response to
something a character in the narration did, or that was done to them.

In \emph{The Grand Tour} (Jones 1974), the narrator as character
expresses regret at having lost a significant other while the song
\emph{He Stopped Loving Her Today}, the narrator, giving a narrative
account of someone else also expresses the regret felt by the character
when their relationship dissolves. The above are emblematic country
songs in that regret is commonly expressed in country music. Different
songs express regret in different contexts however. \emph{Apologize} by
Timbaland is not a country song, and expresses regret in a different
context. Here the narrator expresses regret for the love that he, the
narrator, \emph{had} felt.

Love songs often express regret when a relationship ends. In the song
\emph{The Grand Tour}, our attention is drawn to the daily use of items
``where she used to . . .'' which speak to the fact that they were once
shared by people who experienced positive sentiments toward one another.
Their current state however suggests that they are no longer used this
way ``. . . this old house will never be the same . . .''. In \emph{The
Grand Tour}, there is existing controversy over whether the end of the
relationship is caused by the narrator, or someone else. In the podcast
\emph{Cocaine and Rhinestones: The history of country music}, Tyler
Mahan Coe suggests that the change in the state of the lives of the
individuals involved in \emph{The Grand Tour} is the result of a tragic
mishap citing these lines in the song; ``There's her rings, all her
things, and her clothes are in the closet'' and ``as you leave you'll
see the nursery, . . . taking nothing but the baby and my heart''.

The events in \emph{The Grand Tour} may not have been caused by the
involvement from any agent. For instance, if Coe is correct, then the
dissolution of the relationship depicted in \emph{The Grand Tour} has
nothing to do with any actions committed on the part of the narrator or
his lost love. In \emph{He Stopped Loving Her Today}, the regret is more
definitively caused by an action when the subject whom the song speaks
about is left by their significant other effectively dissolving their
relationship. ``He said I'll love you `till I die, She told him 'you'll
forget in time'\,''. Regret can be the result of either an action or an
event and the regret in \emph{The Grand Tour} is likely caused by an
event, whereas in Jones' other song, \emph{He Stopped Loving Her Today},
it is likely caused by an action for instance if the protagonist pursued
undisclosed intimate relationships with others.

While in the two previous examples \emph{regret} is expressed upon the
end of a relationship or rather its dissolution, in other examples, more
positive emotions, not regret, are expressed upon the end or dissolution
of the relationship rather than of the relationship's continuance. For
instance Kelly Clarkson's \emph{Since U Been Gone} does this: ``Since
you've been gone, I can breath for the first time . . .''. There are
also different reasons for the regret whether expressed for having
suffered through the ordeal of that relationship, or for its end.
Sometimes the regret is expressed as a longing for various qualities the
loved possessed, for instance in the \emph{Grand Tour}, ``over there,
sits the chair \ldots{} and whisper `oh I love you'\,''. For these
reasons, narratives can help us to feel closeness and union with the
loved one even though they are fictional.

Arguably, when a subject experiences this sort of disassociation
identified with suffering above, they are no longer in a position to
experience regret when loosing the things most dear to them. This
correlates with expressions of emotion associated with love, where love
has often been defined as the feeling of loving emotions provoked in one
party by another. Upon the object's absence, the subject no longer feels
these emotions and experiences regret on account of this absence (Taylor
1976). The lover can also experience regret in longing for the qualities
not possessed by the loved, when someone whom they loved failed to
reciprocate in word or deed, or failed in some other way. For instance
in Timbaland's \emph{Too Late to Apologize} the following line expresses
this: ``tell me that you need me then you go and cut me down''.

\section{Identity Reductionism}\label{identity-reductionism}

Narrative accounts of identity attempt to use the way an individual
responds to events in a person's life to identify the agent through
space and time. For instance, in the agency literature, reductionists
explain the metaphysical continuity of an object across space and time
in terms of personal identity. Derek Parfit (1987) argues that in order
to understand what sort of object a person is, we need to understand two
different qualities. These include the nature of a person and their
\emph{criterion of identity}, i.e., their continued existence over time.
But when we consider these qualities in contrast to persons versus
non-persons, we include not only physical states of matter, but also
immaterial ones such as mental states, memories, beliefs, motivations,
etc. A person can be identified just if continuity between subsequent
immaterial mental states and physical states are achieved.

Reductionists such as Parfit can define personal identity as the
persistence of the human person over time, and say that this is
reducible to the binding of psychological states with physical matter.
Causalism is the most popular account that explains the origin of
psychological states, an object \emph{A} at time \emph{T\(_1\)} is
identical to object \emph{B} at time \emph{T\(_2\)} as long as we can
draw a continuous line from object \emph{A} at time \emph{T\(_1\)} to
object \emph{B} at time \emph{T\(_2\)}. The continuous line here is most
often mental states such as memory, passions, character traits and
dispositions. But often times it is the case that those memories we
accept as definitive, have their origins in persons or events other than
ourselves. There are certain dispositions I have because of persons
other than myself signing various documents articulating what rights I
recognize as belonging to me for instance.

While Parfit's account does show how persons are not mere matter, he and
other causalists state that an action performed by an agent, is an event
caused by the mental state of that agent (Altshuler 2021). The
difficulty arises in understanding how mental states are related to
prior events. While a mental state may be the source of some event,
another event may be the source of another mental state. Regardless of
the relationship between mental states and the material universe, it is
not clear where the causal chain originates.

Korsgaard (1989) has argued that on Parfit (1987)`s account of identity,
a person is primarily a \emph{locus of experience} (Korsgaard 1989,
103). This means that the way to identify an object as existing through
space and time as it correlates with a person, is by identifying a set
of experiences between objects across space and time that are connected
through a set of events or some other connecting factor. These events
are psychological states and 'spatiotemporally continuous living
bod{[}ies{]}' (Korsgaard 1989, 104). In this way I might define myself
as \emph{an American} and defend such values as independence and
freedom. It is plausible then to draw a line from events in 1776,
America's founding, to my present day. Therefore I associate various
values I hold with events that transpired in 1776 and as such, the
psychological states I currently have, have their own origins centuries
before. In other words, because a person's psychological states have
prior causes, the states themselves are merely some form of event and
experience.

Narrative accounts then are responses to impersonal accounts. According
to narrative accounts, how we frame events such as these, is what makes
the event meaningful for a given individual (Schechtman 1996, 96).
Without narratives, the agency and the identity of the subject is lost
among a causal chain of events. This has been called `constitution'. In
terms of the constitution model of identity, without the binding effect
of narrative, the subject of a given event is victim to a disparate
causal chain of cause and effect. Our identity then would merely be an
impersonal metaphysical continuation of matter (Korsgaard 1996).

Altshuler identifies a problem with the above account, namely that it
does not account for those events that are not selected by the agent. He
argues that such events do have a causal effects on what ultimately
constitutes one's identity. Following Peter Goldberg (2009), Altshuler
instead identifies narrative with an agent's \emph{ability} to
reconfigure their motivational and psychological past with respect to
their own reactions to these events. In other words, Altshuler too,
defends a self-constitution model of agency and identity, but one on
which a person's responses to his or her past, shape her future
identity.

At the center of the difference between popular narratives like those
found in popular songs and narrative identity accounts is that popular
narratives are merely ones we are prompted to imagine (Gilmore 2020,
132). But the source of narrative identities is not clear. The question
then is what connection do those narratives have to those that we are
prompted to imagine. John Gilmore (2020) suggests that the relationship
is weak and dependent on the function of the imagination for engaging
with the work.

One way narrative identity may relate to popular narratives, is if a
response is required by the artist that maps onto what an audience
perceives as normative. However, Gilmore (2020) rejects this arguing
that different popular narratives have different responses. Further,
these responses are not justified by normative expressions of emotion,
where the emotion elicited by the context is one we would expect outside
the popular narrative {[}p.~132{]}. Rather, Gilmore defends an
intentionalist stance of normative discontinuity. According to this
stance, an author's intentions constrain and justify the emotional
response of their audiences. Accordingly, laughter which would otherwise
be inappropriate would be an apt response to some fiction. It is not
clear whether this view can justify one's identification with a racist
trope of themselves however.

Patricia Hill Collins and Wendy Laybourn on the other hand suggest a
strong relationship (Collins 2006, 2009; Laybourn 2018). The narratives
in popular media are controlling images designed to justify racist and
sexist attitudes.

\begin{quote}
Controlling images ``are designed to make racism, sexism, poverty, and
other forms of social injustice appear to be natural, normal, and
inevitable parts of everyday life'' (Collins 2009, 77).
\end{quote}

It is important to point out here that I am not asking about purpose or
the design of such prompted imaginings. Rather that focusing on
Laybourn's context of `authenticity' and `realness' highlights a popular
implicit assumption regarding what is true and what is fiction in
popular media. Peter Goldberg (2009) highlights our emotions as a form
of reflection, especially that of regret. Sometimes as agents, the
expressions of regret enables one to better exercise their own agency.
We can imagine someone being cheated through the purchase of a used car.
Looking towards the future, that individual would exercise more caution
as a consequence upon reflection of their mistake. In these moments, the
ability to tell a story becomes important as we are enabled by the
narrative to recognize the appropriate objects of our emotions in the
story. Reflecting on the failed used car purchase, the subject
experiences regret. This subsequent experience pushes them to exercise
more caution in the future. Similarly, we can compare this to another
example. Here, regret is either appropriate to aspects of the story
which bear no direct relation on the narrator, or ones that do. Being
able to recognize and reflect on emotions enables one to express their
agency going forward. What regret is afforded to the subject when
reflecting on decisions made by themselves?

In this context, narrative is substantial in our lives as it provides
ways of thinking about how our lives may have gone differently, or may
go differently and why; our plans for the future and how things might
turn out.\footnote{For contrast, see Michael Bratman (2000) argues that
  as agents, we are purposeful and reflective about our motivations, and
  capable of organizing our activity over and in-line with our purposes
  and motivations as we reflect on them. One method of reflection
  involves the expressions of various emotions upon the occurrence of
  some event.} Thinking in this way leads to regret if we are
disappointed in a particular choice we have made. Our feelings of regret
can cause us to behave differently in the future. Questions about
narrative thinking then just are questions about moral agency.

Importantly, the possible prospect of being able to imagine today what a
morally good life will look like for us tomorrow, is a question about
how what we believe about ourselves today, including what desires we
have, influence our ability to pursue these ways of being. These in turn
determine our regrets. Imagining what kind of person our spouse will be
can perhaps open up some romantic possibilities while closing others.
Imagining where we will buy our first home may restrict our career
choices, just like imagining how many children we will have may affect
where we want to live. But more importantly, having various beliefs
about ourselves under gird all these other possibilities.

More importantly, we can develop a sense of agency when we imagine ways
in which our friendships, family ties, and other important relationships
will go. These are preceded by beliefs in us about ways they
\emph{could} go. We can imagine the type of person we are to others and
take needed steps for becoming that person, if we believe it is
possible. As such, the person we will be and how we will relate to
others is the goal that is formed by a narrative. Consider one
particular narrative, we imagine ourselves as thoughtful, kind,
considerate, loving, or compassionate. Imagining ourselves this way
becomes the motivation we need to read up on the latest parenting
methods. We might imagine ourselves being invited to witness the birth
of a friend's first child and scour the internet for advice on how to be
a better gift giver.

While the argument that the way we imagine our lives going can have a
positive causal effect on the kind of person we become, one may argue
that this view is too optimistic. To what extent do the experiences we
have shape the person we become? For instance, Eric J. Miller (2020)
argues that victimhood imposes normative standards that ground the moral
status of the victim. What he means is that social roles, such as
victim, determine the values we are expected to live by.

As a victim, there are various normative standards imposed on us that
undermine our ability to imagine other parenting styles and values, such
as those relating to how we protect and nourish our children.
Libertarian values may be those we are less apt to promote in our homes.
One reason is that when victims as parents, or the parents of victims,
instruct a child, they must instruct them to act in self-subordinating
ways ``when confronted by the police'' {[}221{]} rather than in ways
that assert their autonomy.

Consider the following example narratives:

\begin{itemize}
\tightlist
\item
  Ty Dolla's \emph{by Yourself}: ``Stacked your bread and bought your
  own Mercedes (Vroom, vroom). You your own boss, do it your way (Way)
\end{itemize}

The song's protagonist is faced with a given environment that forces her
to exert her indepedence rather than her social bonds. As such, her lack
of regret is a normative standard imposed on her. We can contrast this
situation with the possibility that while material goods may or may not
be a necessary condition for well-being, they are certainly not a
sufficient one as expressed in other popular songs. The protagonist in
the song \emph{Rich and Sad} by Post Malone, looks back over other
events, motives, and desires that resulted in the dissolution of his
relationship and he feels sadness, as an expression of regret:

\begin{itemize}
\tightlist
\item
  ``All this stuntin' couldn't satisfy my soul (--oul), Got a hundred
  big places, but I'm still alone (--one)''.
\end{itemize}

First, each of the above narratives are constrained by beliefs the
narrator may have about themselves. In the first, the narrator,
believing they have no other choice but to cut themselves off from
anyone who is unable to contribute financially. But it does lead to
another criticism, the agency of the artist versus that of the audience.
Some music criticisms will often point out that musical taste assumes a
sense of moral autonomy not available to some. If an artist is
accustomed to abject poverty, then that person might be seen as
vulnerable to market forces, creating art that is in high demand for the
purposes of financial gain (Laybourn 2018), or only capable of
representing their immediate physical environment (Richardson and Scott
2002). Regret within narratives are an example of such narrative agency.
But if Eric Miller is correct, then some populations just do not possess
agency in the requisite sense. But this view would thoroughly undermine
decades of social justice movements in law and elsewhere.

The common responses to such criticisms come in two forms. They
typically make some causal claim about the effects of environment on
both artistic performance (Richardson and Scott 2002), and audience
music selection criteria, for instance, that artists are merely defining
their personal experiences (Shusterman 2007, 419). This responses
highlight unequal wealth opportunities imposing normative standards on
their victims. The second kind of response makes a cultural claim
deriving from supposed traditional practices in pre-colonial Africa
(Abrahams 2006).

Importantly here, the responses that are grounded in assertions about
the culture of pre-colonial Africa are not necessarily a response to the
empirical claims about the representations of anti-social behavior in
rap and hip-hop music, but rather to the metaphorical use of language
exhibited by pre-colonial African tribal practices. What are commonly
noted as anti-social themes in rap and hip-hop then are downplayed as
word-play underlying a cultural or political critique or something else
entirely. Namely the references to violence and misogyny are
allegorical, not literal (Smitherman 1997; Gates 1989).

Magnus and Malone (2023) give such an account by highlighting the
criminal prosecutions that cite song lyrics as testimonial confessions.
Malone argues that being too generous in treating artistic assertions as
truth in songs, commits the harm of art-interpretive injustice. She
argues that the listener, in this case the prosecution, extends the
wrong degree of credibility to the artist. Instead of interpreting the
song lyrics as art, the listener interprets the lyrics as testimony.

Importantly, a component of Malone's argument is that prosecutions do
not treat other genres of music based on the content of their art.
However, there is another and more important issue at concern here. This
worry is found in the dual character account of art Malone has
previously defended. There, Malone argues that \emph{authentic} country
music is a dual character art-form that embodies ``the core value
commitments of the genre'' (Malone 2023, 1). According to this view,
country music, like punk rock and hip-hop (Liao, Meskin, and Knobe 2020)
constitute core value commitments. As such, following the dual character
account of art, it is not clear what level of credibility we should
assign to the song lyrics in criminal prosecutions.

Consider the following dialogue used by Liao, Meskin, and Knobe (2020),
7-8:

\subsubsection{The Scientist}\label{the-scientist}

\begin{quote}
Alfie: She is not a scientist. I know many people think she is a
scientist, but when you think about what scientists really are, you will
realize that she is not a scientist at all.
\end{quote}

\begin{quote}
Betty: Of course she is a scientist. She is skilled at running
experiments and publishing in scientific journals.
\end{quote}

\begin{quote}
Alfie: That is not what I meant! What I meant is that she does not
possess a genuine willingness to revise her beliefs in light of
empirical evidence.
\end{quote}

\subsubsection{The Christian}\label{the-christian}

\begin{quote}
Alfie: She is not a Christian. I know many people think she is a
Christian, but when you think about what it really means to be a
Christian, you will realize that she is not a Christian at all.
\end{quote}

\begin{quote}
Betty: Of course she is a Christian. She goes to church every week and
believes that Jesus is the son of God.
\end{quote}

\begin{quote}
Alfie: That is not what I meant! What I meant is that she does not show
genuine compassion.
\end{quote}

Malone follows Liao, Meskin, and Knobe (2020) who argue that we can
``pick out the concepts'' for which people make various judgments as
those above, by distinguishing them from the kinds of judgments ``from
concepts for which people would not make such judgments'' (Liao, Meskin,
and Knobe 2020, 6). Here, many things will both satisfy and at the same
time fail to satisfy the criteria that allows one to make identity
judgments. I argue that here, as a character attribution, these concepts
enable one to judge whether someone \emph{truly} belongs to a given
category.

What is important among these responses however, is that they too
undermine attributions of personal identity and human agency and go
against the dual character concept of art. A question they force us to
ask, is what songs belong in a given category? Separately, some have
held that we can determine the values an artist holds, by carefully
listening to the song lyrics they write. Although this fact would
suggest that an artist's lyrics might qualify as a personal narrative
defining an artist's identity and agency, it is not clear whether such
narratives would qualify on this account of narrative agency.

Altshuler's account could give more credence to music lyrics as an
example narrative in the psychological sense, it is not clear how this
may be the case. According to some accounts, an artist's written lyrics
and musical arrangement merely express plausible values someone might
hold (Malone 2023), while others articulate them as more definitive
(Meissner and Huebner 2022), still others argue that art is functional
(Gilmore 2020), thereby leaving open the possibility they could count as
narrative in the functional sense, but do not do so necessarily.

For instance, Whyte (2018), Meissner and Huebner (2022) have argued that
\emph{settler colonial imaginaries} define what they call \emph{settler}
ecologies. These are narratives that replace existing narratives. They
are pro-social concepts that are central for members of a given
community and include beliefs that are central to a group's flourishing.
They can include songs that invoke tropes of indigenous identities. Tim
McGraw's song \emph{Indian Outlaw} is an example of this sort of
practice. In the song he claims to be half Cherokee and half Choctaw.
But it is ambigous whether he is referring to the character, or if the
song is loosely biographical. Either way, it offers a portrait about
what it means to be indigenous American. Importantly, the song is not
aiming at the truth of what it means to be indigenous American. As such,
Meissner and Huebner (2022) would argue that it \emph{fortifies
ignorance} of the people and the land that they inhabit. Most
importantly, it undermines relational networks between the people that
would sustain their \emph{co-flourishing}.

It may be that rap and hip-hop are not good example narratives. But they
show promise. Consider the question above, is some material good a
necessary and sufficient condition for well-being as expressed in some
popular songs? As narratives, the two examples above would better
explain why some event happened. Why the protagonist in \emph{by
yourself} severed their relationship with their significant other, and
why the protagonist in \emph{Rich and Sad} is rich and sad.

Here, we primarily contrast {[}money{]} as the material object
expressing that it:

\begin{itemize}
\tightlist
\item
  ``could satisfy my soul''
\end{itemize}

And contrasting this with {[}money{]} as the material object expressing
that it:

\begin{itemize}
\tightlist
\item
  ``could not satisfy my soul''.
\end{itemize}

One way that narratives are thought to explain agency then, is by
synchronizing incongruous events, states, and motives. As such, a
narrative can say yes, money is either a necessary and sufficient
condition for well-being for some agent or not. As such, the expressions
of regret offers a plausible way to recognize the core of one's agency.

Recognizing the form of some plausible narratives may help to understand
how to recognize a narrative if we say it. The various incidents and
events have made it difficult for the agent to recognize value in
non-material objects. We can, and often do, ask a meta-ethical question
about what kinds of objects are truly of moral value and why. The
narrative provides us with a plausible account of the why. For instance,
we can ask to what extent the quest for consumer goods ought to outweigh
other pursuits such as time with one's children or calling one's mother
on Sunday. What is probably missing from the above is the understanding
that there is a difference between only some material possession being
necessary for well-being, versus some non-material possession being
necessary for well-being.

But exactly how are material and non-material objects expressed over and
against one another other? It may be the case that volunteering at an
orphanage is of moral value and it can be expressed in song in at least
two ways. Lyrically, ``I wish I had volunteered at the orphanage
yesterday''. Or lyrically that expresses a negative outcome for not
having volunteered at the orphanage yesterday. In constructing lyrics of
the second form, the author chooses words that either have been
historically encoded with the sentiment she is hoping to ascribe to a
state of affairs, or is potentially encoded with that sentiment in such
a way that the desired emotion is potentially ascribed to that state of
affairs. However, if I neglect my children in volunteering at an
orphanage, then \emph{my} volunteering is of reduced moral value and we
could imagine a lyrics saying ``I wish I had spent time with my children
instead of volunteering at the orphanage yesterday'', or that says ``I
spent so many hours at the orphanage that I do not know who my own
children are anymore''; or ``I spent hours at the orphanage, I do not
know who my own children are anymore''.

We can also express a normative question about what behaviors are
appropriate, when we should be kind, or whether a given action is more
kind than another using similar methods. But there are additional
interesting questions. As inferred above, we ask about the context in
which a given ethical question is of moral value, and we can express it
as a value by inviting the audience to reflect on such expressions by
coercing their emotions accordingly. Volunteering my time at an
orphanage is of moral value, and so I would use music or poetry to evoke
emotions in the audience that correspond with this value. Considering
the context is a way to assess the rationality of a given question,
emotional expression in art is a way to make sense of that question.
Making sense of the question helps us to anticipate the answer in at
least two ways, what does the artist suggest is valuable, we would then
look at the emotional expressions in her art, and what is the agent
potentially capable of valuing.

For instance, assuming that Jill is the last sentient creature on a far
away planet, it may not make sense to ask what Jill's ethical
obligations are to those around her. But contexts regarding ethical
questions are important. There are two important such contexts here;
socio-economic or cultural environments, and psychological or social
ones. Given questions about morality, our concerns are largely about how
we should behave toward others. For many years, this was answered by
referencing cultural and socio-economic conditions. More
controversially, it has also been answered by referencing psychological
or social conditions. But this is not the same as asking whether a given
question about ethical obligation makes sense or not. It would not make
sense to say a person with a medical condition that precludes the
experience of empathy, has an ethical obligation to experience the
emotions that align with empathy. There is a way in which our questions
about moral behavior should make sense. This is in the context of our
behavior toward others. Most humans, will generally live in contexts
wherein our behavior to others should have a kind of rationality about
them, in valuing our relationships with others.

The social version of this question is more opaque. It is a question
about our social environments and to what extent our social environments
can provide us with the resources necessary to fulfill our moral
obligations. Once again, there are two categories of question here. For
instance, consider our choices as to what kind of careers would make
fulfilling our moral obligations possible. If we recognize a moral
obligation to give vast amounts of wealth to charity, then we could
pursue a set of social connections (work hard in high school and get
into a good university) that would make our ability to fulfill this
obligation more likely. The other category of question is even more
difficult however. It begins as a question about ourselves and the
sources of our beliefs as to what is possible, and what is not. We would
start with a question about what is true about ourselves. Related to
this is a question about what extent do those around us contribute to
our beliefs about what is true about ourselves and what is obligatory in
regard to our behavior. If we consider that we have a moral obligation
to support national health, then there are some political parties we
might campaign on behalf of and others we would not. Here I suggest then
that how we should understand such questions is in the form of a
narrative.

\chapter{Comments on Reconstructions of Ballet
Classics}\label{comments-on-reconstructions-of-ballet-classics}

Montaque Reynolds (Stetson University)

\hfill\break

\begin{itemize}
\item
  How should the contemporary dance community assess the authenticity of
  classical ballet reconstructions? Two important possiblities include
  the work's musical notations, and the work's text or story (Van Camp
  unpublished draft, 1). The distinction here seems to be between a
  work's syntax, akin to its grammatical structure, and its semantics,
  or meaning.
\item
  On the one hand, the authenticity of a musical work is determined by
  its compliance with a musical score (Kania 2024). In ballet, we have
  the musical score and so it is plausible that we determine the work's
  authenticity through a positive assessment of its compliance with a
  musical score and dance score, but for reconstructions of
  nineteenth-century works, we do not have the dance notations. Yet
  despite the lack of a work's dance score, we still consider the work
  authentic.
\item
  Therefore it is not clear whether the authenticity of a ballet
  performance also relies solely on compliance with its musical score.
  Traditionally this has meant that the choreography is amenable to
  change without affecting our understanding of the authenticity of the
  performance.
\item
  Professor Julie Van Camp then asks what factors in 19th century ballet
  reconstructions and reenactments help us to determine whether a given
  dance is authentic, accurate, or neither.
\item
  One possibility is requiring the same sets of standards that would be
  employed when considering a historical item for display in a museum.
\item
  One helpful heuristic of the contrast between reenactment and
  reconstruction then, is the focus on demonstrating how a dance might
  be redone versus a phenomenological way of understanding ``what it was
  \emph{like to do it}'' (Franko 2017, 1). One helpful label of this
  work then would be to focus on the authenticity of a given performance
  either as being compliant with a score, or compliance with the
  original meaning of the work.
\item
  Using the definition offered by Mark Franko (2017) in his introduction
  to the Oxford handbook of dance and reenactment, reconstructions are
  preservations of past accomplishments while it is not yet clear
  whether reenactments aim to reimagine those past performances instead,
  or whether these have some additional aim.
\item
  On this later point, it would appear that there are many important
  reasons for re-imagining a ballet performance.
\item
  For one, if we do not have an authentic historical item from a
  particular time period to include in a museum collection from that
  time period, it may make sense to use the nearest counterpart.
\item
  Analogously, considering the 19th century ballet performances
  \emph{Swan Lake}, \emph{Giselle} and \emph{The Sleeping Beauty}, these
  exist though we do not have relevant dance notations of their original
  performances.
\item
  Therefore it would be an important project to determine what the
  nearest counterpart might be since there are at least two different
  ways to understand the term ``nearest counterpart''. One important
  direction of consideration here, is that an important part of the
  essence of a performance, apart from its compliance with either a
  music or dance score would be the emotional essence of a given
  performance.
\end{itemize}

Lukacs (1977) argues that a work's emotional essence is produced by
social factors.

\begin{itemize}
\item
  But these do not seem to be factors that can be displayed in a museum
  in the traditional sense.
\item
  Perhaps the closest terms we have for this distinction in the
  reenactment and reconstruction literature, is what Andrew Kania (2024)
  attributes to Julian Dodd (2020); `compliance' authenticity versus
  \emph{`interpretive'} authenticity.
\item
  First, Professor Van Camp carefully and expertly draws out the
  complications for understanding why this is a hard problem.
\item
  In three prominent examples of the reconstructions of the 19th century
  pieces, \textbf{\emph{Swan Lake}}, \textbf{\emph{Giselle}}, and
  \textbf{\emph{Sleeping Beauty}}, if we focus on audience reactions to
  these pieces, then the `reconstructions' do not appear to have been as
  appealing as their 19th century reenactment counterparts since
  audience favor typically follows works that reimagine classics.
\item
  This later position is explored with more detail in a review written
  by Professor Van Camp explaining the loose requirements for
  authenticity in ballet reenactments versus those of reconstructions,
  usually focus on a shift of social values.
\item
  Tamara Rojo as the artistic director for the San Francisco Ballet
  delicately introduces variations in her version of \emph{Raymonda} by
  ``expunging Middle Eastern stereotypes'' and exploring new forms of
  agency by introducing a female heroine rather than a male hero.
\item
  This shows the tolerance ballet has for variation without loss of
  ontological identity.
\item
  But this may also just mean that a work's ontological identity is
  grounded somewhere else.
\item
  Professor Van Camp defends these loose standards as authentically
  derivative since often times our own standards are amenable to change
  in other areas, such as values which we know have strong emotive
  force.
\item
  A central question here then could be the shifting plasticity of
  standards and how these diverge between reconstructions and
  reenactments.
\item
  This is especially important if we consider whether there is a more
  ultimate and underlying truth at which multiple performances of a work
  may attempt to reveal.
\item
  At the Paris 1913 premier of \emph{Le Sacre du Printemps} staged by
  Vaslav Nijinsky, it is said that the audience rioted owing to the near
  primal movements of the performers.
\item
  In 2009, the resident choreographer for the Boston Ballet premiered a
  version of \emph{Le Sacre du Printemps} that diverged from Nijinsky's
  ``uncontrolled and raw movements'' with the introduction of ``clean
  shapes, intricate isolations, and dynamic partner work''.

  \begin{itemize}
  \tightlist
  \item
    Elo defends his interpretation arguing that highlighting the
    performances universal value of ``community, identity, and social
    systems'' did not require the ``sacrifice for the common good'' that
    it required when Nijinsky staged the work in 1913.
  \end{itemize}
\item
  This is important considering the role the performance played in
  moving the theater into the modern era of ballet.
\item
  But this also leads to an additional question, how audiences think of
  these 19th century reconstructions as authentic and historically
  accurate and why and what this means for our accepted notions of
  authenticity and accuracy.
\item
  Perhaps it is not necessary here to make any comment on what is
  considered an audience.
\item
  Some audiences desire to `view the past accomplishments of dance' and
  others are more interested in imagining what a historical performance
  would look like when adding the benefit of `improved diet, health, and
  training regimes' of the modern period.
\end{itemize}

As noted by Julian Dodd (2020) in his \emph{Being True to Works of
Music}, it is possible that the distinction here then, is one between
`compliance' authenticity and `interpretive' authenticity.

\begin{itemize}
\item
  For instance, it is not clear whether Professor Van Camp follows
  Andrew Kania (2024) in equating `authenticity' with `genuineness' so
  that an `authentic' performance is a mere `genuine' interpretation of
  a score, and more importantly what this has to do with the compliance
  and interpretation distinction.
\item
  However, one plausible direction in which this discussion could move,
  is through a focus on interpretive authenticity.
\item
  The biggest complication here is that taking the understanding of
  producing a `genuine' work means producing an `authorized' score for
  that work.
\item
  But this is not all that determines the authenticity of a work.
\item
  There are many considerations that go into producing an `authorized'
  score.
\end{itemize}

Lukacs (1977) would argue that these considerations include the
underlying values the score seeks to express.

\begin{itemize}
\item
  Professor Van Camp considers Swan Lake, The Sleeping Beauty, and
  Giselle since there are instances of these works that were created
  without a choreographic notional system.
\item
  The inclusion of the notational system in order to be authorized then,
  needs to account of fidelity to some expressed value.
\item
  Therefore, this fact will have us question the conditions for
  understanding a work rely solely on the \emph{musical} notation
  ascribed to a ballet performance and not specific features of the
  dance itself.
\item
  Therefore there are at least two different standards for a work's
  authenticity.
\item
  So while it is not clear what the standards are for a work being
  authentic, * Professor Van Camp may accept the ballet reconstruction
  as authentic simply if it complies with the musical score.
\item
  But using the 1913 performance of \emph{Le Sacre du Printemps} in
  Paris staged by Nijinsky and the 2009 version staged by Elo, another
  important question emerges.
\item
  Here, maybe there are underlying values which a work intends to
  portray and we can either expect riots when a performance fails to
  express them adequately, or acclaim when they succeed.
\item
  With ballet, which can contain both musical and choreographic scores,
  it would seem that there are additional factors one may consider when
  attempting to determine authenticity besides fidelity to a
  performances musical, or dance score.
\item
  So even if we only consider the musical score for a historic ballet
  performance, there are still multiple types of authenticity we can
  rely on and what the scores these are is not clear.
\end{itemize}

\chapter{Puzzles and Mysteries}\label{puzzles-and-mysteries}

Montaque Reynolds (Stetson University)

\hfill\break

\chapter{Scope, Outline}\label{scope-outline}

\begin{enumerate}
\def\labelenumi{\arabic{enumi}.}
\tightlist
\item
  Use open world/sand box/MMO video games to create unique thought
  experiments through video game play.
\item
  Use these thought experiments to motivate moral reasoning and moral
  decision making.
\item
  Read and analyze classical literature to understand the use of
  literature and multi-media in moral development.
\item
  Conceptually develop games based on scenes in the literature we read.
\end{enumerate}

I recently defended a dissertation in philosophy. One of the questions
that philosophers discuss is the way in which personal identity, who we
understand ourselves to be, is socially and culturally constructed. Some
more well known examples include music, e.g., folk songs, cultural
dances, etc. However, more recently discussions about video games have
begun to become prevalent in this area, for instance consider the unjust
depiction of urban community members in Grand Theft Auto.

Art can be thought of as a form of counterfactual argument, a way to
negotiate a common practical identity.\footnote{Malone (2023), 14} One
form of the representational arts, e.g., music (feelings), stage
performances, movies, etc.

Consider: moral reasons why an emotion is unmerited (Gilmore 2020, 130)

\begin{itemize}
\tightlist
\item
  Laughing at the victim of bullying
\end{itemize}

\begin{enumerate}
\def\labelenumi{\arabic{enumi}.}
\tightlist
\item
  Phillip L. Hammack ``Narrative and the Cultural Psychology of
  Identity'', \emph{process} of identity development represents the link
  between self and society'' (Hammack 2008, 224)
\item
  Marya Schechtman ``Personhood and Personal Identity'', self-knowledge:
  `Is this really what I believe?'; `Is this really what I desire?'; `Is
  this really what I intend?'; ``The resources involved in answering
  these questions would be resources that allowed a person to assign
  particular psychological states to herself . . .'' (Schechtman 1990,
  89)
\item
  Peter Goldie, ``Narrative Thinking, Emotion, and Planning'', ties
  feelings of regret to what he calls `narrative thinking', but might be
  better understood as narrative identity or `narrative agency'.
\end{enumerate}

\begin{quote}
Now, in narrative thinking, when we look back with regrets on what we
did and said, and we look forward in making self-governing policies in
the light of those regrets, we \emph{necessarily} see ourselves as
agents persisting over time. More expanded type of representation
wherein the player does not merely need to follow the reasoning of the
creator of the artwork, but can exercize even more human agency perhaps
exemplifying her own character even more.
\end{quote}

The open-world, MMO, sandbox nature of the developing video game, space
sim Star Citizen is a great example of this. Whereas games like the GTA
series make it obvious and explicit that success in the game can only be
had by exercizing what are typically called vices, Star Citizen lets the
player define success and methods of success. What Robson and Meskin
(2016) may call ``interactive fictions''.

Here, the game's developer relinquishes a modicum of agency so that the
player can develop her own. Not only execution element, but uniquely
agential as well.

In these, the player can follow one of the many narratives together, or
combine several together into their own unique narrative experience. Or
they could choose to ignore to ignore any of the narratives and either
explore, grief others which is a kind of wholly novel narrative
experience, or create their own narratives either in a piece-meal
fashion (allow their choices to define the narrative arc), or design and
build a narrative.

\section{Project Description}\label{project-description}

\begin{enumerate}
\def\labelenumi{\arabic{enumi}.}
\tightlist
\item
  Year 1: Open world games and custom moral thought experiments
\item
  Year 2: Group readings of literary works and writing projects aimed at
  developing video games
\item
  Year 3: Learning how to use a game engine and publish a game
\item
  Year 4: Perfecting the game and learning how to make a portfolio
\item
  Post high school: Mentor current students
\item
  Post high school: Mentor current high school students and 1st year
  university students
\end{enumerate}

It will be a four year after school program. Year one will begin by
inviting them to play open world games, games that do not prescribe any
sort of game play but allows the player to choose their own in-game
objectives.

The next stage (year 2) will involve group readings of literary works
that invite the reader to imagine themselves as a member of some given
culture. In the second semester the students will develop a project
wherein they will come up with a plan for creating their own work of art
that invites the audience to imagine themselves as a member of a given
culture this will be a writing project.

In year 3, the students will spend the first semester learning how to
use a game engine and the second semester they will learn how to publish
their game.

Year 4, the students will perfect their game, and learn how to make a
portfolio.

Ultimately, I would like the students to come back after high school to
mentor current students, whether they are working or during breaks from
university.

\section{Budget}\label{budget}

\subsection{Curriculum Development}\label{curriculum-development}

\begin{itemize}
\tightlist
\item
  Design Curriculum for 1 course per semester for beginning

  \begin{itemize}
  \tightlist
  \item
    Develop literature-based game concepts (e.g., adapting Nicomachean
    Ethics scenes).
  \item
    Hire instructional designers, game developers, and literature
    scholars. (Not confident about this direction)
  \end{itemize}
\item
  Build Learning Management System (LMS) for the course

  \begin{itemize}
  \tightlist
  \item
    This may be done by using grok (already in progress)
  \item
    However, if needed to right away, can temporarily use open source
    LMS, e.g., Frappe (SLU?)

    \begin{itemize}
    \tightlist
    \item
      Free, but requires a hosting plan
    \item
      Best option is Oracle Cloud Free Tier
    \item
      For duration of \$300 credit
    \item
      \$10-\$20/month afterwards
    \item
      Second best is to self-host on linux machine
    \end{itemize}
  \end{itemize}
\end{itemize}

\subsection{Pilot the Program}\label{pilot-the-program}

\begin{itemize}
\tightlist
\item
  Pilot the program with 10 students

  \begin{itemize}
  \tightlist
  \item
    YMCA, Schools, BWorks alumni, libraries
  \item
    Star Citizen Licenses, \$450 per semester x 4
  \item
    Set up guest speakers, Graduate students, interested professors
  \item
    SLU: use esports lab. Offer space to faculty for their students.
  \end{itemize}
\end{itemize}

\subsection{Evaulation and Research}\label{evaulation-and-research}

\begin{itemize}
\tightlist
\item
  Assesss moral reasoning outcomes, tests, papers, etc. Plus qualitative
  feedback, interviews, etc.
\item
  Publish results in peer-reviewed journals
\end{itemize}

\subsection{Seattle Expansion and Scaling (2--3 years):
\$200,000--\$600,000}\label{seattle-expansion-and-scaling-23-years-200000600000}

\begin{itemize}
\tightlist
\item
  Adapt curriculum for Seattle schools and esports organizations (e.g.,
  Seattle Surge).
\item
  Partner with tech-driven institutions (e.g., University of
  Washington).
\item
  Market to game studios (e.g., Bungie).
\end{itemize}

\chapter{Presence in Pedagogy: At the Heart of Learning
Philosophy}\label{presence-in-pedagogy-at-the-heart-of-learning-philosophy}

The perceptual illusion of nonmediation:

\begin{itemize}
\tightlist
\item
  Communication, psychology is an attempt to understand how we
  understand and interpret information in nonmediated settings, store
  and retrieve information, memories, make decisions. (3)
\end{itemize}

\begin{enumerate}
\def\labelenumi{\arabic{enumi}.}
\tightlist
\item
  Presence as social richness: ``presence is the extent to which a
  medium is perceived as sociable, warm, sensitive, personal, or
  intimate when it is used to interact with other people. (Short,
  Williams, and Christie, 1976)'' , (5).
\item
  Presence as realism: ``presence is the degree to which a medium can
  produce seemingly accurate representations of objects, events, and
  people---representations that look, sound, and/or feel like the `real'
  thing'' (6).
\item
  Presence as transportation: ``user is transported to another place''.
\item
  Presence as immersion: ``the degree to which the virtual environment
  submerges the perceptual system of the user'' (Biocca \& Delaney,
  1995, p.~57).
\item
  Presence as social actor within medium: user's responses ``to social
  cues presented by persons they encounter within a medium even though
  it is illogical and even inappropriate to do so''. {[}See Lombard,
  1995 and Lemish, 1982{]}
\item
  Presence as medium as social actor: ``social responses of media users
  to cues provided by the medium itself {[}13{]}.
\end{enumerate}

But what is the distinction of presence in mediated and nonmediated
environments? Arguably, if the illusion is not an illusion, but fact
about the environment, then we would say that the person is present.

\begin{quote}
An ``illlusion of nonmediation'' occus when a person fails to perceive
or acknowledge the existence of a medium in his/her communication
environment and responds as he/she would if the medium were not there.
{[}13{]}
\end{quote}

Nonmediated: experienced without human-made technology

Nonmediated presence, experience w/o human-made tech Mediated presence,
experience \emph{as though} w/o human-made tech

\chapter{Simulating Philosophy and other Lab
Experiments}\label{simulating-philosophy-and-other-lab-experiments}

\section{Simulating Philosophy by Marcus
Schulzke}\label{simulating-philosophy-by-marcus-schulzke}

\begin{itemize}
\tightlist
\item
  Schulzke proposes studying video games as executable narrative thought
  experiments.
\end{itemize}

Thought experiments in philosophy refer to hypothetical scenarios that
are used to challenge or substantiate a theory.

James (1993): narratives that lead readers or listeners through a
sequence of events that reflect how one would experience the real world.
In this way, they are abstract models of more complex issues. They
elemenate or simplify extraneous information in order to focus on the
elements of an event taht raise a theorectical challenge.

\begin{itemize}
\tightlist
\item
  Used to challenge a theory by identifying a difficult case a theory
  cannot answer.
\item
  Can model situations that may be untestable using empirical methods.
\item
  Less dangerous than real world experiments.

  \begin{itemize}
  \tightlist
  \item
    For instance, consider experiments about euthanasia or abortion.
  \item
    Duties to one another (One hour one life)
  \item
    Cooperative activity (Eve Online, Star Citizen)
  \end{itemize}
\end{itemize}

\section{What is Going On?}\label{what-is-going-on}

\subsection{Thought Experiment One}\label{thought-experiment-one}

Nicholas Rescher (1991):

\begin{itemize}
\tightlist
\item
  Thought experiments can be used to assess theorectical and empirical
  research questions.
\end{itemize}

\begin{quote}
A `thought experiment' is an attempt to draw instruction from a process
of hypothetical reasoning that proceeds by eliciting the consequences of
an hypothesis which, for aught that one actually knows to the contrary,
may well be false. It consists in reasoning from a supposition that is
not accepted as true---perhaps is even known to be false---but is
assumed provisionally in the interests of making a point or resolving a
conclusion.
\end{quote}

\begin{itemize}
\tightlist
\item
  It is to be shown that \emph{P} is the case (where it has not yet been
  established whether \emph{P} or not \emph{-P}).
\item
  Assume---as ``thought experiment''---that \emph{P} is the case (which
  is not inherently implausible).
\item
  Explain \emph{Q} on the basis of this assumption, where \emph{Q} is
  something patently true which we could not readily explain otherwise.
\item
  Hence maintain that \emph{P}.
\end{itemize}

\subsection{Thought Experiment Two}\label{thought-experiment-two}

John Norton (1991):

Thought experiments are arguments which:

\begin{enumerate}
\def\labelenumi{\arabic{enumi}.}
\tightlist
\item
  Posit hypothetical or counterfactual states of affairs, and
\item
  Invoke particulars irrelevant to the generanlity of the conclusion.
\end{enumerate}

\begin{itemize}
\tightlist
\item
  \emph{hypothetical or counterfactual} opposite of actual.
\item
  \emph{irrelevant particulars} irrelevant to the conclusion.
\end{itemize}

\begin{quote}
Thus, in one version of the thought experiement in which Einsteinn
sought to demonstrate that the effects of acceleration mimic those of
gravitation, he asked us to imagine a physicist-observer who has been
drugged and reawakens closed up inside a box (Einstein, 1912,
pp.~1254-55). . . . Without particulars such as these, however, thought
experiments would not have their experimental appearance.
\end{quote}

Good thought experiments are good arguments. Either:

\begin{enumerate}
\def\labelenumi{\arabic{enumi}.}
\tightlist
\item
  The particular-free conclusion follows deductively form the premisses.
  Thought experiments of this type are typically \emph{reductio}
  arguments.
\item
  The conclusion is freed of the particulars by an inductive step. This
  step might involve the assertion that the case involving the
  particulars is `typical' or that the particulars are `inessential,' so
  that the result derived holds in other cases as well.
\end{enumerate}

\part{Drafts \& Notes}

\chapter{Notes}\label{notes}

There are several ways in which some art objects may represent more
complex realities. One may think that the purpose of a painting is to
reflect the experiences one had during a war, or peace time event. Such
a painting may express the given facts of the case, for instance an
artist could paint the signing of the Magna Carta.

If the painting were painted by an artist who was alive to witness the
signing of the Magna Carta, then the artist would merely be representing
contemporary reality. If the artist were not alive at the time, then we
could argue that the artist is representing historical reality. The
artist is relying on trusted and well known sources would be different
from the artist imagining an event in which not many direct sources
exist. But an artist can represent a projection of reality as well when
the reality they want to convey is obscure, not well known, highly
contested or hypothetical.

Moral (2021) argues that art's purpose in the middle ages is to
represent reality. It is one thing to accurately and factually represent
a tree, however, Pinero is talking about more complex reality,
especially moral reality. One example of medieval representations of
moral reality include the Persian Manticore. The manticore as depicted
in medieval art is said to represent evil.

According the this definition of medieval art, art implies metaphysical
aspirations and normative elements (Also see De Bruyne 1959). Its
commitment to reality is undeniable, yet the reality which it attempts
to represent is so complex and vast that the limits between what is and
what is not, between what we believe and what we know, between what we
are and what we wish to be, are often blurred (\textbf{cron41?}).

Here, the artist can represent facts as they believe them to be, or as
they wish them to be. It is this latter motivation that Pinero suggests
motivated medieval artwork representing the manticore. There is a
difference here then between modern and medieval artistic function. This
in itself is uninteresting. However, what may be more interesting to the
reader, is the further claim that this difference can be applied in
response to the contemporary account of epistemic injustice.

\subsection{Wendy Laybourn's Cost of Being
Real}\label{wendy-laybourns-cost-of-being-real-2}

Wendy Laybourn (2018) draws a comparison between black authenticity and
what she calls commodified realness. According to Laybourn, identity is
sold to a broader public as black authenticity. This criticism is
centered around the effects of what she calls \emph{controlling images}
in popular media. Of deep concern here, is that many of the members of
the communities depicted by such images, consider aspects of these
images to be authentic representations of themselves and the communities
in which they live. Therefore, there is a public perception that such
images are not in themselves harmful. If they are in some degree
controlling, then they would be harmful, but there is strong
disagreement regarding whether or not they are (20C History Project
2014; Collins 2006, 2009; Sartwell 1998).

\chapter{Currie's Fiction}\label{curries-fiction-1}

Why do we think that fairy tales are fiction but news stories are not?
Is fake news the same thing as a fairy tale? During the pandemic, fake
news became a widely discussed talking point. A less widely discussed
talking point however was why don't we just call it fiction? A
derivative of this question is how fiction is made possible. Meaning,
what are the creative differences between fake news, news news and fairy
tales? To answer this question, Gregory Currie (1990) suggests three
possible answers.

According to the front page news, the terrorist leader Amelia Boyd was
killed by the Headhunters Gang. First are the characteristics that
distinguish fiction from non-fiction. This question is a close
investigation of the entities that are in fiction and make them
possible. For instance, the psychological and linguistic resources that
producers and consumers of fiction must have access to.

First, consider that fiction uses language. But there is no universal
linguistic feature shared by all fictional works and there is no
available proof/s confirming what is a fiction. We usually think that
reading a work is an act of reading words and sentences that make it
up.2 However, consider a mime? Or silent films? These are works that do
not use language in the traditional sense, yet we consider them
fictional. Therefore, what makes something a fiction cannot necessarily
be a linguistic feature. Therefore what makes something a fiction is
neither a necessary nor sufficient condition for being fictional.
Finally then, something can be a fiction because of its verbal
structure, but what makes it a fiction does not lie in the text itself.

\section{Semantic Properties}\label{semantic-properties-1}

Currie concludes that fictionality is a relational property, it is
something possessed in virtue of the text's relations to other things.
Therefore while it is possible that fictionality is a linguistic matter,
what makes it a fiction lies in some other property. (a method of
communication or other communicative practice)

On this account, fictional properties, like non-fictional ones, have the
semantic properties of truth value and reference. What this means is
that a fictional statement can be true and false. For instance, consider
the statement ``Killing an NPC in monitored space will earn you a crime
state''. According to Currie, this statement is true. What makes
something true or false in a fiction, is its relation to the outside
world and so what is important then, is what this relationship is.

Because fiction is a relational property, they are community based and
``depend on the prevailing attitudes of the community, or of a subgroup
of the community rather than the ideas of a single individual'' (Currie
1990, 10).

\begin{quote}
For I take it that an act of joint authorship is exactly that: an act
engaged in by more than one person rather than several distinct acts
undertaken individually and patched together. (Currie 1990, 11)
\end{quote}

The process of creating a fiction begins with the intention. It is not a
singular intention however, but is a social enterprise. So maybe the
reason that we do not call fake news, fiction, is because of the
purposes and intent of fake news. It is intended to deceive rather than
express true propositions. But if this is true, then why do we not
consider that fiction is intended to deceive? Why does fiction have
semantic properties such as truth whereas fake news does not? Currie
(1990)

Owing to the fantastical nature of these accounts, it is highly probable
that they are either misguided representations of some more familiar
animal, or motivated by purposes other than truth. In the classical era,
it is most likely the case that the Manticore is a confused
representation of some other more familiar animal, while in the medieval
period, the motivation for the representation lies outside of truth.

On this second point, in the 57th question of the 1st of the second part
of the Summa Theologica, Thomas Aquinas defends the view of art as a
form of right reason. There is a distinction here between the right
reason for making objects, and the right reason for doing good actions.
We should first focus on a central point here, what is important to both
of these activities, namely the right use of one's cognitive faculties.
Therefore, the difference here lies in the aesthetic representation,
whether of an action or object.

\phantomsection\label{refs}
\begin{CSLReferences}{1}{0}
\bibitem[\citeproctext]{ref-20c14}
20C History Project. 2014. {``Video {Music Box}: {Rap Debate} with
{Ice-T} and {Calvin Butts}.''}

\bibitem[\citeproctext]{ref-abra06}
Abrahams, Roger D. 2006. \emph{Deep down in the {Jungle}: {Black
American Folklore} from the {Streets} of {Philadelphia}}. New Brunswick,
N.J: Transaction Publishers.

\bibitem[\citeproctext]{ref-alts21}
Altshuler, Roman. 2021. {``Agency, {Narrative}, and {Mortality}.''} In
\emph{The {Routledge Handbook} of {Philosophy} of {Agency}}, 1st ed.,
385--93. London: Routledge.
\url{https://doi.org/10.4324/9780429202131-43}.

\bibitem[\citeproctext]{ref-brat00}
Bratman, Michael E. 2000. {``Reflection, {Planning}, and {Temporally
Extended Agency}.''} \emph{The Philosophical Review} 109 (1): 35--61.
\url{https://doi.org/10.2307/2693554}.

\bibitem[\citeproctext]{ref-coe24}
Coe, Tyler Mahan. 2024. \emph{Cocaine and {Rhinestones}: {A History} of
{George Jones} and {Tammy Wynette}}. 1st ed. New York: Simon \&
Schuster.

\bibitem[\citeproctext]{ref-coll06}
Collins, Patricia Hill. 2006. \emph{Black {Sexual Politics}: {African
Americans}, {Gender}, and the {New Racism}}. Sociology {Race} \&
{Ethnicity}. New York: Routledge.

\bibitem[\citeproctext]{ref-coll09}
---------. 2009. \emph{Black {Feminist Thought}: {Knowledge},
{Consciousness}, and the {Politics} of {Empowerment}}. 2nd ed. Routledge
{Classics}. New York: Routledge.
\url{https://doi.org/10.4324/9780203309506}.

\bibitem[\citeproctext]{ref-curr90}
Currie, Gregory. 1990. {``Supervenience, {Essentialism} and {Aesthetic
Properties}.''} \emph{Philosophical Studies} 58 (3): 243--57.

\bibitem[\citeproctext]{ref-debr59}
De Bruyne, Edgar. 1959. \emph{Estudios de Est{é}tica Medieval}. Vol. 17.
Editorial Gredos.

\bibitem[\citeproctext]{ref-dodd20}
Dodd, Julian. 2020. \emph{Being {True} to {Works} of {Music}}. Oxford,
GB: Oxford University Press.

\bibitem[\citeproctext]{ref-fran17}
Franko, Mark. 2017. \emph{The {Oxford} Handbook of Dance and
Reenactment}. Oxford Handbooks. New York (N.Y.): Oxford university
press.

\bibitem[\citeproctext]{ref-fric07}
Fricker, Miranda. 2007. \emph{Epistemic {Injustice}: {Power} and the
{Ethics} of {Knowing}}. Oxford ; New York: Oxford University Press.

\bibitem[\citeproctext]{ref-gate89}
Gates, Henry Louis. 1989. \emph{The {Signifying Monkey}: {A Theory} of
{African-American Literary Criticism}}. Oxford University Press, USA.

\bibitem[\citeproctext]{ref-gilm20}
Gilmore, Jonathan. 2020. \emph{Apt {Imaginings}: {Feelings} for
{Fictions} and {Other Creatures} of the {Mind}}. 1st ed. Oxford
University Press.
\url{https://doi.org/10.1093/oso/9780190096342.001.0001}.

\bibitem[\citeproctext]{ref-gold09}
Goldberg, Sanford. 2009. {``Experts, {Semantic} and {Epistemic}.''}
\emph{Nous} 43 (4): 581--98.
\url{https://doi.org/10.1111/j.1468-0068.2009.00720.x}.

\bibitem[\citeproctext]{ref-hamm08}
Hammack, Phillip L. 2008. {``Narrative and the {Cultural Psychology} of
{Identity}.''} \emph{Pers Soc Psychol Rev} 12 (3): 222--47.
\url{https://doi.org/10.1177/1088868308316892}.

\bibitem[\citeproctext]{ref-harr11}
Harrison, Mary-Catherine. 2011. {``How {Narrative Relationships Overcome
Empathic Bias}: {Elizabeth Gaskell}'s {Empathy} Across {Social
Difference}.''} \emph{Poetics Today} 32 (2): 255--88.
\url{https://doi.org/10.1215/03335372-1162686}.

\bibitem[\citeproctext]{ref-jack20}
Jackson, Zakiyyah Iman. 2020. \emph{Becoming {Human}: {Matter} and
{Meaning} in an {Antiblack World}}. Vol. 53. NYU Press.
\url{https://www.jstor.org/stable/j.ctv1n6ptnn}.

\bibitem[\citeproctext]{ref-jone74}
Jones, George. 1974. {``The {Grand Tour}.''} Nashville TN.

\bibitem[\citeproctext]{ref-jone79}
---------. 1979. {``He {Stopped Loving Her Today}.''} Nashville TN: Epic
Records.

\bibitem[\citeproctext]{ref-kani24}
Kania, Andrew. 2024. {``The {Philosophy} of {Music}.''} In \emph{The
{Stanford Encyclopedia} of {Philosophy}}, edited by Edward N. Zalta and
Uri Nodelman, Winter 2024. Metaphysics Research Lab, Stanford
University.

\bibitem[\citeproctext]{ref-kors89}
Korsgaard, Christine M. 1989. {``Personal {Identity} and the {Unity} of
{Agency}: {A Kantian Response} to {Parfit}.''} \emph{Philosophy \&
Public Affairs} 18 (2): 101--32.
\url{https://www.jstor.org/stable/2265447}.

\bibitem[\citeproctext]{ref-kors96}
---------. 1996. \emph{The {Sources} of {Normativity}}. Cambridge:
Cambridge University Press.

\bibitem[\citeproctext]{ref-layb18}
Laybourn, Wendy M. 2018. {``The Cost of Being {`Real'}: Black
Authenticity, Colourism, and {Billboard Rap Chart} Rankings.''}
\emph{Ethnic and Racial Studies} 41 (11): 2085--2103.
\url{https://doi.org/10.1080/01419870.2017.1343484}.

\bibitem[\citeproctext]{ref-liao20}
Liao, Shen-yi, Aaron Meskin, and Joshua Knobe. 2020. {``Dual {Character
Art Concepts}.''} \emph{Pacific Philosophical Quarterly} 101 (1):
102--28. \url{https://doi.org/10.1111/papq.12301}.

\bibitem[\citeproctext]{ref-luka77}
Lukacs, Georg, ed. 1977. {``{Realism in the Balance}.''} In
\emph{{Aesthetics and politics}}. London: NLB.

\bibitem[\citeproctext]{ref-magn23}
Magnus, P. D., and Evan Malone. 2023. {``Popular {Music} and
{Art-Interpretive Injustice}.''} \emph{Inquiry: An Interdisciplinary
Journal of Philosophy}.
\url{https://doi.org/10.1080/0020174x.2023.2213740}.

\bibitem[\citeproctext]{ref-malo23}
Malone, Evan. 2023. {``Country {Music} and the {Problem} of
{Authenticity}.''} \emph{Brit J Aesthetics} 63 (1): 75--90.
\url{https://doi.org/10.1093/aesthj/ayac020}.

\bibitem[\citeproctext]{ref-mbem17}
Mbembe, Achille. 2017. \emph{Critique of {Black Reason}}. Translated by
Laurent Dubois. Durham: Duke University Press.

\bibitem[\citeproctext]{ref-meis22}
Meissner, Shelbi Nahwilet, and Bryce Huebner. 2022. {``Outlaw
{Epistemologies}: {Resisting} the {Viciousness} of {Country Music}'s
{Settler Ignorance}.''} \emph{Philosophical Issues} 32 (1): 214--32.
\url{https://doi.org/10.1111/phis.12229}.

\bibitem[\citeproctext]{ref-mill20}
Miller, Eric J. 2020. {``The {Moral Burdens} of {Police Wrongdoing}.''}
\emph{Res Phil.} 97 (2): 219--69.
\url{https://doi.org/10.11612/resphil.1901}.

\bibitem[\citeproctext]{ref-mora21}
Moral, Ricardo Piñero. 2021. {``Aesthetics of {Evil} in {Middle Ages}:
{Beasts} as {Symbol} of the {Devil}.''} \emph{Religions} 12 (11).
\url{https://doi.org/10.3390/rel12110957}.

\bibitem[\citeproctext]{ref-morr19}
Morrison, Elizabeth, and Larisa Grollemond. 2019. \emph{Book of Beasts:
The Bestiary in the Medieval World {Exhibition}, {Los Angeles}, {J}.
{Paul Getty Museum}, from {May} 14 to {August} 18, 2019}. Los Angeles:
J. Paul Getty Museum.

\bibitem[\citeproctext]{ref-nguy14}
Nguyen, Jenny, and Amanda Koontz Anthony. 2014. {``Black {Authenticity}:
{Defining} the {Ideals} and {Expectations} in the {Construction} of
'{Real}' {Blackness}.''} \emph{Sociology Compass} 8 (6): 770--79.
\url{https://doi.org/10.1111/soc4.12171}.

\bibitem[\citeproctext]{ref-parf87}
Parfit, Derek. 1987. \emph{Reasons and {Persons}}. 1. issued in
paperback (with corr.), reprinted with further corr. Oxford: Clarendon
Press.

\bibitem[\citeproctext]{ref-rich02}
Richardson, Jeanita W., and Kim A. Scott. 2002. {``Rap {Music} and {Its
Violent Progeny}: {America}'s {Culture} of {Violence} in {Context}.''}
\emph{The Journal of Negro Education} 71 (3): 175.
\url{https://doi.org/10.2307/3211235}.

\bibitem[\citeproctext]{ref-robs16}
Robson, Jon, and Aaron Meskin. 2016. {``Video {Games} as {Self-Involving
Interactive Fictions}.''} \emph{The Journal of Aesthetics and Art
Criticism} 74 (2): 165--77. \url{https://doi.org/10.1111/jaac.12269}.

\bibitem[\citeproctext]{ref-sart98b}
Sartwell, Crispin. 1998. \emph{Act Like {You Know}: {African-American
Autobiography} \& {White Identity}}. Chicago, Ill.: Univ. of Chicago
Press.

\bibitem[\citeproctext]{ref-sche90}
Schechtman, Marya. 1990. {``Personhood and {Personal Identity}.''}
\emph{The Journal of Philosophy} 87 (2): 71.
\url{https://doi.org/10.2307/2026882}.

\bibitem[\citeproctext]{ref-sche96}
---------. 1996. \emph{The {Constitution} of {Selves}}. Cornell
University Press.

\bibitem[\citeproctext]{ref-shus07}
Shusterman, Richard. 2007. {``Rap as {Art} and {Philosophy}.''} In,
419--28. \url{https://doi.org/10.1002/9780470751640.ch29}.

\bibitem[\citeproctext]{ref-smit97b}
Smitherman, Geneva. 1997. {``"{The Chain Remain} the {Same}":
{Communicative Practices} in the {Hip Hop Nation}.''} \emph{Journal of
Black Studies} 28 (1): 3--25.
\url{https://www.jstor.org/stable/2784891}.

\bibitem[\citeproctext]{ref-stum10}
Stump, Eleonore. 2010. \emph{Wandering in {Darkness}: {Narrative} and
the {Problem} of {Suffering}}.

\bibitem[\citeproctext]{ref-tayl76}
Taylor, Gabriele. 1976. {``{VIII}*---{Love}.''} \emph{Proceedings of the
Aristotelian Society} 76 (1): 147--64.
\url{https://doi.org/10.1093/aristotelian/76.1.147}.

\bibitem[\citeproctext]{ref-vanunpublisheddraft}
Van Camp, Julie. unpublished draft. {``{RECONSTRUCTIONS OF BALLET
CLASSICS Accepted}.''}

\bibitem[\citeproctext]{ref-whyt18}
Whyte, Kyle. 2018. {``Settler {Colonialism}, {Ecology}, and
{Environmental Injustice}.''} \emph{Environment and Society} 9: 125--44.
\url{https://www.jstor.org/stable/26879582}.

\end{CSLReferences}




\end{document}
